\section{AMC 12 12A 2018}
        \begin{enumerate}
        	\item (\textit{AMC 12 12A 2018}) A large urn contains $100$ balls, of which $36 \%$ are red and the rest are blue. How many of the blue balls must be removed so that the percentage of red balls in the urn will be $72 \%$? (No red balls are to be removed.)


 $\textbf{(A)}\ 28 \qquad\textbf{(B)}\  32 \qquad\textbf{(C)}\  36 \qquad\textbf{(D)}\   50 \qquad\textbf{(E)}\ 64$

 \textbf{Answer: }D



\textbf{Solution 1}

There are $36$ red balls; for these red balls to comprise $72 \%$ of the urn, there must be only $14$ blue balls. Since there are currently $64$ blue balls, this means we must remove $\boxed{\textbf{(D)}\ 50}$.

 



\textbf{Solution 2}

There are $36$ red balls and $64$ blue balls. For the percentage of the red balls to double from $36 \%$ to $72 \%$ of the urn, half of the total number of balls must be removed. Therefore, the number of blue balls that need to be removed is $\boxed{\textbf{(D)}\ 50}$.

 



\textbf{Solution 3}

There are $36$ red balls out of the total $100$ balls. We want to continuously remove blue balls until the percentage of red balls in the urn is 72\%.Therefore, we want
 \[\frac{36}{100-x}=\frac{72}{100}.\]Solving for $x$ gives that we must remove $\boxed{\textbf{(D)}\ 50}$ blue balls.

 


	\item (\textit{AMC 12 12A 2018}) While exploring a cave, Carl comes across a collection of $5$-pound rocks worth $\$14$ each, $4$-pound rocks worth $\$11$ each, and $1$-pound rocks worth $\$2$ each. There are at least $20$ of each size. He can carry at most $18$ pounds. What is the maximum value, in dollars, of the rocks he can carry out of the cave?


 $\textbf{(A) } 48 \qquad \textbf{(B) } 49 \qquad \textbf{(C) } 50 \qquad \textbf{(D) } 51 \qquad \textbf{(E) } 52$

 \textbf{Answer: }C



\textbf{Solution 1}

The value of $5$-pound rocks is $\$14\div5=\$2.80$ per pound, and the value of $4$-pound rocks is $\$11\div4=\$2.75$ per pound. Clearly, Carl should not carry more than three $1$-pound rocks. Otherwise, he can replace some $1$-pound rocks with some heavier rocks, preserving the weight but increasing the total value. 

 We perform casework on the number of $1$-pound rocks Carl can carry:
 \[\begin{array}{c|c|c||c} & & & \\ [-2.5ex] \boldsymbol{1}\textbf{-Pound Rocks} & \boldsymbol{4}\textbf{-Pound Rocks} & \boldsymbol{5}\textbf{-Pound Rocks} & \textbf{Total Value} \\ \textbf{(}\boldsymbol{\$2}\textbf{ Each)} & \textbf{(}\boldsymbol{\$11}\textbf{ Each)} & \textbf{(}\boldsymbol{\$14}\textbf{ Each)} & \\ [0.5ex] \hline & & & \\ [-2ex] 0 & 2 & 2 & \$50 \\  & & & \\ [-2.25ex] 1 & 3 & 1 & \$49 \\ & & & \\ [-2.25ex] 2 & 4 & 0 & \$48 \\ & & & \\ [-2.25ex] 3 & 0 & 3 & \$48 \end{array}\]Clearly, the maximum value of the rocks Carl can carry is $\boxed{\textbf{(C) } 50}$ dollars.

 <u>\textbf{Remark}</u>

 Note that an upper bound of the total value is $\$2.80\cdot18=\$50.40,$ from which we can eliminate choices $\textbf{(D)}$ and $\textbf{(E)}.$




\textbf{Solution 2}

Since each rock is worth $1$ dollar less than $3$ times its weight (in pounds), the answer is just $3\cdot 18=54$ minus the minimum number of rocks we need to make $18$ pounds. Note that we need at least $4$ rocks (two $5$-pound rocks and two $4$-pound rocks) to make $18$ pounds, so the answer is $54-4=\boxed{\textbf{(C) } 50}.$



	\item (\textit{AMC 12 12A 2018}) How many ways can a student schedule $3$ mathematics courses -- algebra, geometry, and number theory -- in a $6$-period day if no two mathematics courses can be taken in consecutive periods? (What courses the student takes during the other $3$ periods is of no concern here.)


 $\textbf{(A) }3\qquad\textbf{(B) }6\qquad\textbf{(C) }12\qquad\textbf{(D) }18\qquad\textbf{(E) }24$

 \textbf{Answer: }E



\textbf{Solution 1}

We must place the classes into the periods such that no two classes are in the same period or in consecutive periods.

 Ignoring distinguishability, we can thus list out the ways that three periods can be chosen for the classes when periods cannot be consecutive:

 Periods $1, 3, 5$

 
Periods $1, 3, 6$

 Periods $1, 4, 6$

 Periods $2, 4, 6$

 There are $4$ ways to place $3$ non distinguishable classes into $6$ periods such that no two classes are in consecutive periods. For each of these ways, there are $3! = 6$ orderings of the classes among themselves.

 Therefore, there are $4 \cdot 6 = \boxed{\textbf{(E) } 24}$ ways to choose the classes.

 -Versailles15625

 



\textbf{Solution 2}

Counting what we don't want is another slick way to solve this problem. Use PIE (Principle of Inclusion and Exclusion) to count two cases: 1. Two classes consecutive, 2. Three classes consecutive. 

 Case 1: Consider two consecutive periods as a "block" of which there are 5 places to put in(1,2; 2,3; 3,4; 4,5; 5,6). Then we simply need to place two classes within the block, $3 \cdot 2$. Finally we have 4 periods remaining to place the final math class. Thus there are $5 \cdot 3 \cdot 2 \cdot 4$ ways to place two consecutive math classes with disregard to the third. 

 Case 2: Now consider three consecutive periods as a "block" of which there are now 4 places to put in(1,2,3; 2,3,4; 3,4,5; 4,5,6). We now need to arrange the math classes in the block, $3 \cdot 2 \cdot 1$. Thus there are $4 \cdot 3 \cdot 2 \cdot 1$ ways to place all three consecutive math classes. 

 By PIE we subtract Case 1 by Case 2 in order to not overcount:  $120-24$. Then we subtract that answer from the total ways to place the classes with no restrictions: $(6 \cdot 5 \cdot 4) - 96= \boxed{\textbf{(E) } 24}$

 -LitJamal

 



\textbf{Solution 3}

We can tackle this problem with a stars-and-bars-ish approach. First, letting math class be 1 and non-math-class be 0, place 0s in between 3 1s:
 \[10101\]Now we need to place 1 additional 0. There are 4 places to put it: 
 \[\underline{\hspace{0.3cm}}1\underline{\hspace{0.3cm}}01\underline{\hspace{0.3cm}}01\underline{\hspace{0.3cm}}\] It can be placed in any 1 of the underscores.Since there are $3!=6$ ways to order the math classes, the answer is $6\cdot 4=\boxed{\textbf{(E) } 24}$.

 



\textbf{Solution 4}

Let $M$ represent a math class and $N$ represent a non-math class. We have \_N\_N\_N\_, where the spaces represent the possible spots the $M’s$ could go in. There are $4\choose 3$ ways to choose the spots for the math classes and $3!=6$ ways to order the classes. Hence, the answer is $4 \cdot 6 = \boxed{\textbf{(E) } 24}$.

 


	\item (\textit{AMC 12 12A 2018}) Alice, Bob, and Charlie were on a hike and were wondering how far away the nearest town was. When Alice said, "We are at least $6$ miles away," Bob replied, "We are at most $5$ miles away." Charlie then remarked, "Actually the nearest town is at most $4$ miles away." It turned out that none of the three statements were true. Let $d$ be the distance in miles to the nearest town. Which of the following intervals is the set of all possible values of $d$?


 $\textbf{(A) }   (0,4)   \qquad        \textbf{(B) }   (4,5)   \qquad    \textbf{(C) }   (4,6)   \qquad   \textbf{(D) }  (5,6)  \qquad  \textbf{(E) }   (5,\infty)$

 \textbf{Answer: }D



\textbf{Solution 1}

For each of the false statements, we identify its corresponding true statement. Note that:

 
            \begin{enumerate} 
        \item $\mathrm{False}\cap\mathrm{True}=\varnothing.$

        \item $\mathrm{False}\cup\mathrm{True}=[0,\infty).$

            \end{enumerate}
            We construct the following table:
 \[\begin{array}{c||c|c} & &  \\ [-2.5ex] \textbf{Hiker} & \textbf{False Statement} & \textbf{True Statement} \\ [0.5ex] \hline & & \\ [-2ex] \textbf{Alice} & [6,\infty) & [0,6) \\  & & \\ [-2.25ex] \textbf{Bob} & [0,5] & (5,\infty)  \\ & & \\ [-2.25ex] \textbf{Charlie} & [0,4] & (4,\infty) \end{array}\]Taking the intersection of the true statements, we have 
 \[[0,6)\cap(5,\infty)\cap(4,\infty)=(5,6)\cap(4,\infty)=\boxed{\textbf{(D) } (5,6)}.\]~MRENTHUSIASM

 



\textbf{Solution 2}

Think of the distances as if they are on a number line. Alice claims that $d > 6$, Bob says $d < 5$, while Charlie thinks $d < 4$. This means that all possible numbers less than $5$ and greater than $6$ are included. However, since the three statements are actually false, the distance to the nearest town is one of the numbers not covered. Therefore, the answer is $\boxed{\textbf{(D) } (5,6)}$.

 


	\item (\textit{AMC 12 12A 2018}) What is the sum of all possible values of $k$ for which the polynomials $x^2 - 3x + 2$ and $x^2 - 5x + k$ have a root in common?


 $\textbf{(A) }3 \qquad\textbf{(B) }4 \qquad\textbf{(C) }5 \qquad\textbf{(D) }6 \qquad\textbf{(E) }10 \qquad$

 \textbf{Answer: }E



\textbf{Solution}

We factor $x^2-3x+2$ into $(x-1)(x-2)$. Thus, either $1$ or $2$ is a root of $x^2-5x+k$. If $1$ is a root, then $1^2-5\cdot1+k=0$, so $k=4$. If $2$ is a root, then $2^2-5\cdot2+k=0$, so $k=6$. The sum of all possible values of $k$ is $\boxed{\textbf{(E) }10}$.

 


	\item (\textit{AMC 12 12A 2018}) For positive integers $m$ and $n$ such that $m+10<n+1$, both the mean and the median of the set $\{m, m+4, m+10, n+1, n+2, 2n\}$ are equal to $n$. What is $m+n$?


 $\textbf{(A) }20\qquad\textbf{(B) }21\qquad\textbf{(C) }22\qquad\textbf{(D) }23\qquad\textbf{(E) }24$

 \textbf{Answer: }B



\textbf{Solution 1}

The mean and median are 
 \[\frac{3m+4n+17}{6}=\frac{m+n+11}{2}=n,\]so $3m+17=2n$ and $m+11=n$. Solving this gives $\left(m,n\right)=\left(5,16\right)$ for $m+n=\boxed{\textbf{(B)}~21}$. (trumpeter)

 



\textbf{Solution 2}

This is an alternate solution if you don't want to solve using algebra. First, notice that the median $n$ is the average of $m+10$ and $n+1$. Therefore, $n=m+11$, so the answer is $m+n=2m+11$, which must be odd. This leaves two remaining options: ${(B) 21}$ and ${(D) 23}$. Notice that if the answer is $(B)$, then $m$ is odd, while $m$ is even if the answer is $(D)$. Since the average of the set is an integer $n$, the sum of the terms must be even. $4+10+1+2+2n$ is odd by definition, so we know that $3m+2n$ must also be odd, thus with a few simple calculations $m$ is odd. Because all other answers have been eliminated, $(B)$ is the only possibility left. Therefore, $m+n=\boxed{21}$. ∎ --anna0kear

 



\textbf{Solution 3}

Since the median is $n$, then $\frac{m+10+n+1}{2} = n \Rightarrow m+11 = n$, or $m= n-11$. Plug this in for $m$ values to get $\frac{7n-16}{6} = n \Rightarrow 7n-16 = 6n \Rightarrow n= 16$. Plug it back in to get $m = 5$, thus $16 + 5 = \boxed{\text{(B)}~21}$.

 


	\item (\textit{AMC 12 12A 2018}) For how many (not necessarily positive) integer values of $n$ is the value of $4000\cdot \left(\frac{2}{5}\right)^n$ an integer?


 $\textbf{(A) }3 \qquad \textbf{(B) }4 \qquad \textbf{(C) }6 \qquad \textbf{(D) }8 \qquad \textbf{(E) }9 \qquad$

 \textbf{Answer: }E



\textbf{Solution 1 (Algebra)}

Note that 
 \[4000\cdot \left(\frac{2}{5}\right)^n=\left(2^5\cdot5^3\right)\cdot \left(2\cdot5^{-1}\right)^n=2^{5+n}\cdot5^{3-n}.\]Since this expression is an integer, we need: 

 
            \begin{enumerate} 
        \item $5+n\geq0,$ from which $n\geq-5.$

        \item $3-n\geq0,$ from which $n\leq3.$

            \end{enumerate}
            Taking the intersection gives $-5\leq n\leq3.$ So, there are $3-(-5)+1=\boxed{\textbf{(E) }9}$ integer values of $n.$




\textbf{Solution 2 (Observations)}

Note that $4000\cdot \left(\frac{2}{5}\right)^n$ will be an integer if the denominator is a factor of $4000$. We also know that the denominator will always be a power of $5$ for positive values and a power of $2$ for all negative values. So we can proceed to divide $4000$ by $5^n$ for each increasing positive value of $n$ until we get a non-factor of $4000$ and also divide $4000$ by $2^{-n}$ for each decreasing negative value of $n$. For positive values we get $n= 1, 2, 3$ and for negative values we get $n= -1, -2, -3, -4, -5$. Also keep in mind that the expression will be an integer for $n=0$, which gives us a total of $\boxed{\textbf{(E) }9}$ for $n.$




\textbf{Solution 3 (Brute Force)}

The values for $n$ are $-5, -4, -3, -2, -1, 0, 1, 2,$ and $3.$

 The corresponding values for $4000\cdot \left(\frac{2}{5}\right)^n$ are $390625, 156250, 62500, 25000, 10000, 4000, 1600, 640,$ and $256,$ respectively.

 In total, there are $\boxed{\textbf{(E) }9}$ values for $n.$



	\item (\textit{AMC 12 12A 2018}) All of the triangles in the diagram below are similar to isosceles triangle $ABC$, in which $AB=AC$. Each of the $7$ smallest triangles has area $1,$ and $\triangle ABC$ has area $40$. What is the area of trapezoid $DBCE$?


 
                    \begin{figure}[H]
                        \centering
                        \includegraphics[width=0.5\textwidth]{Problem/AMC_12_12A_2018_8_Problem_Figure_1.png}
                    \end{figure}
                    
$\textbf{(A) }   16   \qquad        \textbf{(B) }   18   \qquad    \textbf{(C) }   20   \qquad   \textbf{(D) }  22 \qquad  \textbf{(E) }   24$

 \textbf{Answer: }E



\textbf{Solution 1}

Let $x$ be the area of $ADE$. Note that $x$ is comprised of the $7$ small isosceles triangles and a triangle similar to $ADE$ with side length ratio $3:4$ (so an area ratio of $9:16$). Thus, we have 
 \[x=7+\dfrac{9}{16}x.\] This gives $x=16$, so the area of $DBCE=40-x=\boxed{(E) 24}$.

 



\textbf{Solution 2}

Let the base length of the small triangle be $x$. Then, there is a triangle $ADE$ encompassing the 7 small triangles and sharing the top angle with a base length of $4x$. Because the area is proportional to the square of the side, let the base $BC$ be $\sqrt{40}x$. The ratio of the area of triangle $ADE$ to triangle $ABC$ is $\left(\frac{4x}{\sqrt{40}x}\right)^2 = \frac{16}{40}$. The problem says the area of triangle $ABC$ is $40$, so the area of triangle $ADE$ is $16$. So the area of trapezoid $DBCE$ is $40 - 16 = \boxed{24}$.

 



\textbf{Solution 3}

Notice $\big[DBCE\big]=\big[ABC\big]-\big[ADE\big]$.Let the base of the small triangles of area 1 be $x$, then the base length of $\Delta ADE=4x$. Notice, $\left(\frac{DE}{BC}\right)^2=\frac{1}{40}\implies \frac{x}{BC}=\frac{1}{\sqrt{40}}$, then $4x=\frac{4BC}{\sqrt{40}}\implies \big[ADE\big]=\left(\frac{4}{\sqrt{40}}\right)^2\cdot \big[ABC\big]=\frac{2}{5}\big[ABC\big]$Thus, $\big[DBCE\big]=\big[ABC\big]-\big[ADE\big]=\big[ABC\big]\left(1-\frac{2}{5}\right)=\frac{3}{5}\cdot 40=\boxed{24}.$

                    \begin{figure}[H]
                        \centering
                        \includegraphics[width=0.5\textwidth]{Solution/AMC_12_12A_2018_8_Solution_Figure_1.png}
                    \end{figure}
                    




\textbf{Solution 4}

The area of $ADE$ is 16 times the area of the small triangle, as they are similar and their side ratio is $4:1$. Therefore the area of the trapezoid is $40-16=\boxed{24}$.

 



\textbf{Solution 5}

You can see that we can create a "stack" of 5 triangles congruent to the 7 small triangles shown here, arranged in a row above those 7, whose total area would be 5. Similarly, we can create another row of 3, and finally 1 more at the top, as follows. We know this cumulative area will be $7+5+3+1=16$, so to find the area of such trapezoid $BCED$, we just take $40-16=\boxed{24}$, like so. 

 



\textbf{Solution 6}

The combined area of the small triangles is $7$, and from the fact that each small triangle has an area of $1$, we can deduce that the larger triangle above has an area of $9$ (as the sides of the triangles are in a proportion of $\frac{1}{3}$, so will their areas have a proportion that is the square of the proportion of their sides, or $\frac {1}{9}$). Thus, the combined area of the top triangle and the trapezoid immediately below is $7 + 9 = 16$. The area of trapezoid $BCED$ is thus the area of triangle $ABC-16 =\boxed{24}$. 

 



\textbf{Solution 7}

You can assume for the base of one of the smaller triangles to be $\frac{1}{a}$ and the height to be $2a$, giving an area of 1. The larger triangle above the 7 smaller ones then has base $\frac{3}{a}$ and height $6a$, giving it an area of $9$. Then the area of triangle $ADE$ is $16$ and $40-16=\boxed{24}$.

 



\textbf{Solution 8}

You can construct another trapezoid directly above the one shown, with it's bottom length as the top length of the original. Its area would then be 9/16 of the original. Repeating this process infinitely gives us the sequence $7\cdot\left(1+\left(\frac{9}{16}\right)+\left(\frac{9}{16}\right)^2+\left(\frac{9}{16}\right)^3\dots\right)$ . Using the infinite geometric series sum formula gives us $7\cdot\left(\frac{1}{1-\frac{9}{16}}\right)=7\cdot\frac{16}{7}=16$. The triangle's area would thus be $40-16=\boxed{24}$.  

 -ConcaveTriangle

 



\textbf{Solution 9 (weird)}

Note that the area of an isosceles triangle is equivalent to the square of its height. Using this information, the height of the smallest isosceles triangle is $1$, and thus its base is $2.$ 

 Let $h$ be the height of the top triangle. We can set up a height-to-base similarity ratio, using the top triangle and $\triangle{ADE}$. The top triangle has a base of $3\cdot{2}=6$, and $DE=4\cdot{2}=8.$ The height of $\triangle{ADE}$ is $h+1$, therefore our ratio is $\frac{h}{6}=\frac{h+1}{8}$, which yields $h=3$ as our answer. 

 To find the area of the trapezoid, we can take the area of $\triangle{ABC}$ and subtract the area of $\triangle{ADE},$ whose base is $8$ and height $3+1=4$. It follows that the area of $\triangle{ADE}=16$, and subtracting this from $40$ gives us $40-16=\boxed{\text{D) }24}$.

 -Benedict T (countmath1)

 



\textbf{Solution 10 (Pretty standard)}

Since each of the smallest triangles have area $1$, we have the combined area of them is $7$. Also, we know that every triangle is similar to big one. Therefore the triangle above the stack of triangles has area of $3^2 = 9$, because the side length is $3$, so we can just square it to find the area. Now to find the area of trapezoid $BCDE$ we can just do: 
 \[40-(9+7) = 40-16=\boxed{\text{E) }24}\]-jb2015007

 


	\item (\textit{AMC 12 12A 2018}) Which of the following describes the largest subset of values of $y$ within the closed interval $[0,\pi]$ for which

 \[\sin(x+y)\leq \sin(x)+\sin(y)\]for every $x$ between $0$ and $\pi$, inclusive?


 $\textbf{(A) } y=0 \qquad \textbf{(B) } 0\leq y\leq \frac{\pi}{4} \qquad \textbf{(C) } 0\leq y\leq \frac{\pi}{2} \qquad \textbf{(D) } 0\leq y\leq \frac{3\pi}{4} \qquad \textbf{(E) } 0\leq y\leq \pi$

 \textbf{Answer: }E



\textbf{Solution 1}

On the interval $[0, \pi]$ sine is nonnegative; thus $\sin(x + y) = \sin x \cos y + \sin y \cos x \le \sin x + \sin y$ for all $x, y \in [0, \pi]$ and equality only occurs when $\cos x = \cos y = 1$, which is cosine's maximum value. The answer is $\boxed{\textbf{(E) } 0\le y\le \pi}$. (CantonMathGuy)

 



\textbf{Solution 2}

Expanding, 
 \[\cos y \sin x + \cos x \sin y \le \sin x + \sin y\] Let $\sin x =a \ge 0$, $\sin y = b \ge 0$. We have that 
 \[(\cos y)a+(\cos x)b \le a+b\] Comparing coefficients of $a$ and $b$ gives a clear solution: both $\cos y$ and $\cos x$ are less than or equal to one, so the coefficients of $a$ and $b$ on the left are less than on the right. Since $a, b \ge 0$, that means that this equality is always satisfied over this interval, or $\boxed{\textbf{(E) } 0\le y\le \pi}$.

 



\textbf{Solution 3}

If we plug in $\pi$, we can see that $\sin(x+\pi) \le \sin(x)$. Note that since $\sin(x)$ is always nonnegative, $\sin(x+\pi)$ is always nonpositive. So, the inequality holds true when $y=\pi$. The only interval that contains $\pi$ in the answer choices is $\boxed{\textbf{(E) } 0\le y\le \pi}$.

 



\textbf{Solution 4 (Answer Choices)}

Notice that for $y=0$ there is equality and nowhere else because the angle is not being rotated enough. We can just try all the subset of values from the answer choices to try out preferably the largest value for each subset that works as they've special trig values and are the upper bound for the set with values of the $x$ interval.

 From this we notice $E$ has the largest bound from the answer choices that works. Hence our answer is $\boxed{E}$.~Batmanstark

 



\textbf{Solution 5}

Rewrite it as 
 \[\cos y \sin x + \cos x \sin y \le \sin x + \sin y\] then move the terms such that it appears as this 
 \[\sin x (\cos y -1) \le \sin y (1- \cos x)\] We know that on the interval of $[0, \pi]$, $0\le (1-\cos x)\le 2$, and $0\le \sin y$, such that the whole thing on the right hand side is always greater than or equal to $0$. On the left hand side, $0\le \sin x$ and we want $(\cos y -1)\le 0$, which $\boxed{0\le y\le \pi}$ satisfies. ~azure123456

 



\textbf{Solution 6}

Using \href{https://artofproblemsolving.com/wiki/index.php?title=Sum-to-product\_identities}{Sum-to-product identities}, the inequality simplifies to $\cos\left(\frac{x-y}{2}\right)\ge\cos\left(\frac{x+y}{2}\right)$. (This manipulation requires dividing both sides by $\sin\left(\frac{x+y}{2}\right)$, but this value is always positive since $\frac{x+y}{2}\le\pi$ for all given values, so we do not have to flip the sign of the inequality.) This can be interpreted as choosing a value $\frac{x}{2}\in\left[0,\frac{\pi}{2}\right]$, making this value the new center of the cosine graph in the $x$-direction, and taking both points $\frac{y}{2}\in\left[0,\frac{\pi}{2}\right]$ distance from the center. By looking at the graph, we find that for all $x$-values, the change in $y$ always makes it so that the value left of the center is greater than the value right of the center, so our answer is $\boxed{\textbf{(E) } 0\le y\le \pi}$.

 


	\item (\textit{AMC 12 12A 2018}) How many ordered pairs of real numbers $(x,y)$ satisfy the following system of equations?

 \begin{align*} x+3y&=3 \\ \big||x|-|y|\big|&=1 \end{align*}

$\textbf{(A) } 1 \qquad  \textbf{(B) } 2 \qquad  \textbf{(C) } 3 \qquad  \textbf{(D) } 4 \qquad  \textbf{(E) } 8$

 \textbf{Answer: }C



\textbf{Solution 1}

We can solve this by graphing the equations. The second equation looks challenging to graph, but start by graphing it in the first quadrant only (which is easy since the inner absolute value signs can be ignored), then simply reflect that graph into the other quadrants.

 The graph looks something like this:
                    \begin{figure}[H]
                        \centering
                        \includegraphics[width=0.5\textwidth]{Solution/AMC_12_12A_2018_10_Solution_Figure_1.png}
                    \end{figure}
                    Now, it becomes clear that there are $\boxed{\textbf{(C) } 3}$ intersection points.

 (not the author)The lines in the first quadrant are $y=x+1$ and $y=x-1$ and we reflect them into all quadrants as all instances of variables are in a context where flipping the sign does not affect the overall expression.

 



\textbf{Solution 2}

$x+3y=3$ can be rewritten to $x=3-3y$. Substituting $3-3y$ for $x$ in the second equation will give $||3-3y|-y|=1$. Splitting this question into casework for the ranges of $y$ will give us the total number of solutions. 

 $\textbf{Case 1:}$ $y>1$:$3-3y$ will be negative so $|3-3y| = 3y-3.$$|3y-3-y| = |2y-3| = 1$
\begin{verbatim}
Subcase 1: $y>\frac{3}{2}$
\end{verbatim}

$2y-3$ is positive so $2y-3 = 1$ and $y = 2$ and $x = 3-3(2) = -3$
\begin{verbatim}
Subcase 2: $1<y<\frac{3}{2}$
\end{verbatim}

$2y-3$ is negative so $|2y-3| = 3-2y = 1$. $2y = 2$ and so there are no solutions ($y$ can't equal to $1$)

 $\textbf{Case 2:}$ $y = 1$:It is fairly clear that $x = 0.$
$\textbf{Case 3:}$ $y<1$:$3-3y$ will be positive so $|3-3y-y| = |3-4y| = 1$
\begin{verbatim}
Subcase 1: $y>\frac{3}{4}$
\end{verbatim}

$3-4y$ will be negative so $4y-3 = 1$ $\rightarrow$ $4y = 4$. We already have this solution from Case 2 as $y = 1$.

 \begin{verbatim}
Subcase 2: $y<\frac{3}{4}$
\end{verbatim}

$3-4y$ will be positive so $3-4y = 1$ $\rightarrow$ $4y = 2$. $y = \frac{1}{2}$ and $x = \frac{3}{2}$.Thus, the solutions are: $(-3,2), (0,1), \left(\frac{3}{2},\frac{1}{2} \right)$, and the answer is $\boxed{\textbf{(C) } 3}$.

 



\textbf{Solution 3 (Straightforward Casework)}

Note that $||x| - |y||$ can take on one of four values where $x$ and $y$ are: $x + y$ (positive, positive), $x - y$ (positive, negative), $-x + y$ (negative, positive), $-x -y$ (negative, negative). So we have 4 cases which we solve by cancellation with the two equations from the problem:

 \textbf{Case 1: ||x| - |y|| = x+y  }
 \[x+3y=3\]
 \[x+y=1\]

 Subtracting:

 $2y=2 \Rightarrow y=1$ and $x=0$.

 $\text{Result: } (0,1)$

 
\textbf{Case 2: ||x| - |y|| = x-y}
 \[x+3y=3\]
 \[x-y=1\]

 Subtacting:

 $4y=2 \Rightarrow y=\dfrac{1}{2}$ an $x=\dfrac{3}{2}$.

 $\text{Result: } \left(\dfrac{3}{2},\dfrac{1}{2}\right)$

 
\textbf{Case 3: ||x| - |y|| = -x+y}
 \[x+3y=3\]
 \[-x+y=1\]

 Adding: 

 $4y=4 \Rightarrow y=1$ and $x=0$.

 
$\text{Result: } (0,1)$. Since this is the same solution as we got in Case 1, we can not count this (otherwise, we would have overcounted).

 
\textbf{Case 4: ||x| - |y|| = -x-y}
 \[x+3y=3\]
 \[-x-y=1\]

 Adding: 

 $2y=4 \Rightarrow y=2$ and $x=-3$.

 
$\text{Result: } (-3,2)$.

 
\textbf{Answer}

 We solved each case by elimination (either adding the two equations or subtracting), to obtain three solutions: $(0, 1)$, $(-3,2)$, $\left(\dfrac{3}{2}, \dfrac{1}{2}\right)$.

 Our answer is $\boxed{\textbf{(C) } 3}$.

 
\textbf{Contributors}

 



\textbf{Solution 4}

Just as in solution $2$, we derive the equation $||3-3y|-|y||=1$. If we remove the absolute values, the equation collapses into four different possible values. $3-2y$, $3-4y$, $2y-3$, and $4y-3$, each equal to either $1$ or $-1$. Remember that if $P-Q=a$, then $Q-P=-a$. Because we have already taken $1$ and $-1$ into account, we can eliminate one of the conjugates of each pair, namely $3-2y$ and $2y-3$, and $3-4y$ and $4y-3$. Find the values of $y$ when $3-2y=1$, $3-2y=-1$, $3-4y=1$ and $3-4y=-1$. We see that $3-2y=1$ and $3-4y=-1$ give us the same value for $y$, so the answer is $\boxed{\textbf{(C) } 3}$




\textbf{Solution 5}

Just as in solution $2$, we derive the equation $x=3-3y$. Squaring both sides in the second equation gives $x^2+y^2-2|xy|=1$. Putting $x=3-3y$ and doing a little calculation gives $10y^2-18y+9-2|3y-3y^2|=1$. From here we know that $3y-3y^2$ is either positive or negative.

 When positive, we get $2y^2-3y+1=0$ and then, $y=1/2$ or $y=1$.When negative, we get $y^2-3y+2=0$ and then, $y=2$ or $y=1$. Clearly, there are $3$ different pairs of values and that gives us $\boxed{\textbf{(C) } 3}$



	\item (\textit{AMC 12 12A 2018}) A paper triangle with sides of lengths $3,4,$ and $5$ inches, as shown, is folded so that point $A$ falls on point $B$. What is the length in inches of the crease?

                    \begin{figure}[H]
                        \centering
                        \includegraphics[width=0.5\textwidth]{Problem/AMC_12_12A_2018_11_Problem_Figure_1.png}
                    \end{figure}
                    
$\textbf{(A) }   1+\frac12 \sqrt2   \qquad        \textbf{(B) }   \sqrt3   \qquad    \textbf{(C) }   \frac74   \qquad   \textbf{(D) }  \frac{15}{8} \qquad  \textbf{(E) }   2$

 \textbf{Answer: }D



\textbf{Solution 1 (Most Practical)}


                    \begin{figure}[H]
                        \centering
                        \includegraphics[width=0.5\textwidth]{Solution/AMC_12_12A_2018_11_Solution_Figure_1.png}
                    \end{figure}
                    

 First, we need to realize that the crease line is just the perpendicular bisector of side $AB$, the hypotenuse of right triangle $\triangle ABC$. Call the midpoint of $AB$ point $D$. Draw this line and call the intersection point with $AC$ as $E$. Now, $\triangle ACB$ is similar to $\triangle ADE$ by $AA$ similarity. Setting up the ratios, we find that
 \[\frac{BC}{AC}=\frac{DE}{AD} \Rightarrow \frac{3}{4}=\frac{DE}{\frac{5}{2}} \Rightarrow DE=\frac{15}{8}.\]Thus, our answer is $\boxed{\textbf{(D) } \frac{15}{8}}$.

 \textbf{Note} In general, whenever we are asked to make a crease, think about that crease as a line of reflection over which the diagram is reflected. This is why the crease must be the perpendicular bisector of $AB$, because $A$ must be reflected onto $B$.

 



\textbf{Solution 2 (Very speedy, recommended if you are almost out of time)}

Use the ruler and regular paper you brought to quickly draw a 3-4-5 triangle of any scale (don't trust the diagram in the booklet). Very carefully fold the acute vertices together and make a crease. Measure the crease with the ruler. If you were reasonably careful, you should see that it measures somewhat more than $\frac{7}{4}$ units and somewhat less than $2$ units. The only answer choice in that range is $\boxed{\textbf{(D) } \frac{15}{8}}$.

 This is pretty much a cop-out, but it's allowed in the rules technically. This is basically useless for problems without diagrams.

 



\textbf{Solution 3}

Since $\triangle ABC$ is a right triangle, we can see that the slope of line $AB$ is $\frac{BC}{AC}$ = $\frac{3}{4}$. We know that if we fold $\triangle ABC$ so that point $A$ meets point $B$ the crease line will be perpendicular to $AB$ and we also know that the slopes of perpendicular lines are negative reciprocals of one another. Then, we can see that the slope of our crease line is $-\frac{4}{3}$.
                    \begin{figure}[H]
                        \centering
                        \includegraphics[width=0.5\textwidth]{Solution/AMC_12_12A_2018_11_Solution_Figure_2.png}
                    \end{figure}
                    

 Let us call the midpoint of $AB$ point $D$, the midpoint of $AC$ point $E$, and the crease line $DF$. We know that $DE$ is parallel to $BC$ and that $DE$'s length is $\frac{BC}{2}=\frac{3}{2}$. Using our slope calculation from earlier, we can see that$-\frac{DE}{EF}=-\frac{\frac{3}{2}}{EF}=-\frac{4}{3}$. With this information, we can solve for $EF$: 
 \[-4EF=(-\frac{3}{2})(3) \Rightarrow -4EF=-\frac{9}{2} \Rightarrow 4EF=\frac{9}{2} \Rightarrow EF=\frac{9}{8}.\]We can then use the Pythagorean Theorem to find $DF$.
 \[\frac{3}{2}^2+\frac{9}{8}^2=DF^2 \Rightarrow \frac{9}{4}+\frac{9\cdot9}{8\cdot8}=DF^2 \Rightarrow \frac{9\cdot2\cdot8}{4\cdot2\cdot8}+\frac{9\cdot9}{8\cdot8}=DF^2 \Rightarrow \frac{9\cdot2\cdot8+9\cdot9}{8\cdot8}=DF^2 \Rightarrow \frac{9(2\cdot8+9)}{8\cdot8}=DF^2\]
 \[\Rightarrow DF=\sqrt{\frac{9(2\cdot8+9)}{8\cdot8}} \Rightarrow DF=\frac{3\cdot5}{8} \Rightarrow DF=\frac{15}{8}\]Thus, our answer is $\boxed{\textbf{(D) } \frac{15}{8}}$.

 



\textbf{Solution 4}

Make use of the diagram in Solution 3. It can be deduced that $AF=BF$. Let $DF=x$. In $\triangle ADF$, $AF^2=x^2+2.5^2 \Rightarrow AF=\sqrt{x^2+2.5^2}$. Then $FC$ also would be $4-\sqrt{x^2+2.5^2}$.

 In $\triangle BCF$, $BF^2=FC^2+BC^2 \Rightarrow (\sqrt{x^2+2.5^2})^2=(4-\sqrt{x^2+2.5^2})^2+3^2$. After some quick math, we get $\sqrt{x^2+2.5^2}=\frac{25}{8}$.Solving for $x$ will give $x=\frac{15}{8}$.

 $\therefore$ $DF=x=$ $\boxed{\textbf{(D) } \frac{15}{8}}$.

 



\textbf{Solution 5}

Like in Solution 3, we can make use of coordinate geometry to solve this problem. Because the length of the crease is constant no matter the positioning of the triangle, reorient the triangle so it has C at the origin and B and A at $(0,3)$ and $(4,0)$ respectively.

 
                    \begin{figure}[H]
                        \centering
                        \includegraphics[width=0.5\textwidth]{Solution/AMC_12_12A_2018_11_Solution_Figure_3.png}
                    \end{figure}
                    

 We know that each point in the crease is equidistant from $A$ and $B$, so the crease must pass through the midpoint of $AB$, which is $({2,\frac{3}{2}})$ , and be perpendicular to hypotenuse $AB$. The crease therefore has a slope of $\frac{4}{3}$.

 Plugging into point-slope we find that the equation of the crease is, $y-y_1=m(x-x_1)\implies y-\frac{3}{2}=\frac{4}{3}(x-2)\implies y=\frac{4}{3}x-\frac{7}{6}$

 Using 0 for y, we see that the crease intersects $AC$ at $(\frac{7}{8},0)$. 

 By the Distance Formula,

 $(2-\frac{7}{8})^2+({\frac{3}{2}})^2=d^2 \implies d^2=\frac{225}{64} \implies d=\boxed{\textbf{(D) }\frac{15}{8}}$




\textbf{Solution 6}

Let $D$ be a point on AC such that $AD=DB$ and $\triangle{ADB}$ is isosceles. The altitude of this triangle is the crease of the paper triangle.

 Setting $DC=x$, we have a right triangle $\triangle{BCD}$ with legs of x and 3, and, because $AD=DB,$ hypotenuse of $(4-x)$.

 $x^2+3^2=(4-x)^2 \implies x^2+9=16-8x+x^2 \implies 8x=7 \implies x=\frac{7}{8}$

 The area of $\triangle{BCD}  \text{ is thus} = \frac{7}{8}*3*\frac{1}{2} \implies \frac{21}{16}.$ Because $[ADB]=[ABC]-[BCD],$ we can plug in and see that $[ADB]=6-\frac{21}{16}=\frac{75}{16}.$

 Since $A=\frac{bh}{2}$, we get $\frac{75}{16}=\frac{5h}{2} \implies 5h=\frac{75}{8} \implies h=\boxed{\textbf{(D) }\frac{15}{8}}$



	\item (\textit{AMC 12 12A 2018}) Let $S$ be a set of $6$ integers taken from $\{1,2,\dots,12\}$ with the property that if $a$ and $b$ are elements of $S$ with $a<b$, then $b$ is not a multiple of $a$. What is the least possible value of an element in $S$?


 $\textbf{(A)}\ 2\qquad\textbf{(B)}\ 3\qquad\textbf{(C)}\ 4\qquad\textbf{(D)}\ 5\qquad\textbf{(E)}\ 7$

 \textbf{Answer: }C



\textbf{Solution 1}

We start with $2$ because $1$ is not an answer choice. We would have to include every odd number except $1$ to fill out the set, but then $3$ and $9$ would violate the rule, so that won't work.

 Experimentation with $3$ shows it's likewise impossible. You can include $7,11,$ and either $5$ or $10$ (which are always safe). But after adding either $4$ or $8$ we have no more valid numbers.

 Finally, starting with $4,$ we find that the sequence $4,5,6,7,9,11$ works, giving us $\boxed{\textbf{(C)}\ 4}.$




\textbf{Solution 2}

We know that all odd numbers except $1,$ namely $3, 5, 7, 9, 11,$ can be used.

 Now we have $7$ possibilities to choose from for the last number (out of $1, 2, 4, 6, 8, 10, 12$). We can eliminate $1, 2, 10,$ and $12,$ and we have $4, 6, 8$ to choose from. However, $9$ is a multiple of $3.$ Now we have to take out either $3$ or $9$ from the list. If we take out $9,$ none of the numbers would work, but if we take out $3,$ we get 
 \[4, 5, 6, 7, 9, 11.\]The least number is $4,$ so the answer is $\boxed{\textbf{(C)}\ 4}.$




\textbf{Solution 3}

We can get the multiples for the numbers in the original set with multiples in the same original set
 \begin{align*} 1&: \ \text{all elements of }\{1,2,\dots,12\} \\ 2&: \ 4,6,8,10,12 \\ 3&: \ 6,9,12 \\ 4&: \ 8,12 \\ 5&: \ 10 \\ 6&: \ 12 \end{align*}
It will be safe to start with $5$ or $6$ since they have the smallest number of multiples as listed above, but since the question asks for the least, it will be better to try others.

 Trying $4,$ we can get $4,5,6,7,9,11.$ So $4$ works.Trying $3,$ it won't work, so the least is $4.$ This means the answer is $\boxed{\textbf{(C)}\ 4}.$




\textbf{Solution 4}

We partition $\{1,2,\ldots,12\}$ into six nonempty subsets such that for every subset, each element is a multiple of all elements less than or equal to itself: 
 \[\{1,2,4,8\}, \ \{3,6,12\}, \ \{5,10\}, \ \{7\}, \ \{9\}, \ \{11\}.\]Clearly, $S$ must contain exactly one element from each subset:

 
            \begin{itemize}
        \item For $\{5,10\},$ we can select either $5$ or $10.$
            \end{itemize}
            
            \begin{itemize}
        \item For $\{3,6,12\},$ we can select either $6$ or $12.$ Recall that since $9\in S,$ we cannot select $3.$
            \end{itemize}
            
            \begin{itemize}
        \item For $\{1,2,4,8\},$ we can select either $4$ (provided that $12\not\in S$) or $8.$ Recall that since either $6\in S$ or $12\in S,$ we can select neither $1$ nor $2.$
            \end{itemize}
            If $4\in S,$ then the possibilities of $S$ are $\{4,5,6,7,9,11\}$ or $\{4,6,7,9,10,11\}.$ So, the least possible value of an element in $S$ is $\boxed{\textbf{(C)}\ 4}.$
<u>\textbf{Remark}</u>

 There exist multiple such partitions of $\{1,2,\ldots,12\}$ into six nonempty subsets, one of which is 
 \[\{1,2,6,12\}, \ \{3,9\}, \ \{4,8\}, \ \{5,10\}, \ \{7\}, \ \{11\}.\]Regardless of which partition we use, we will conclude that to minimize the least element of $S,$ the only possibilities for $S$ are $\{4,5,6,7,9,11\}$ or $\{4,6,7,9,10,11\}.$




\textbf{Solution 5}

We start with 2 as 1 is not an answer option. Our set would be $\{2,3,5,7,11\}$. We realize we cannot add 12 to the set because 12 is a multiple of 3. Our set only has 5 elements, so starting with 2 won't work. 

 We try 3 next. Our set becomes $\{3,4,5,7,11\}$. We run into the same issue as before so starting with 3 won't work.

 We then try 4. Our set becomes $\{4,5,6,7,9,11\}$. We see we have 6 elements with none being multiples of each other. Therefore our answer is $\boxed{\textbf{(C)}\ 4}.$



	\item (\textit{AMC 12 12A 2018}) How many nonnegative integers can be written in the form 
 \[a_7\cdot3^7+a_6\cdot3^6+a_5\cdot3^5+a_4\cdot3^4+a_3\cdot3^3+a_2\cdot3^2+a_1\cdot3^1+a_0\cdot3^0,\]
where $a_i\in \{-1,0,1\}$ for $0\le i \le 7$?


 $\textbf{(A) } 512 \qquad \textbf{(B) } 729 \qquad \textbf{(C) } 1094 \qquad \textbf{(D) } 3281 \qquad  \textbf{(E) } 59,048$

 \textbf{Answer: }D



\textbf{Solution 1}

This looks like balanced ternary, in which all the integers with absolute values less than $\frac{3^n}{2}$ are represented in $n$ digits. There are 8 digits. Plugging in 8 into the formula for the balanced ternary gives a maximum bound of $|x|=3280.5$, which means there are 3280 positive integers, 0, and 3280 negative integers. Since we want all nonnegative integers, there are $3280+1=\boxed{3281}$ integers or $\boxed{\textbf{D}}$.

 



\textbf{Solution 2}

Note that all numbers formed from this sum are either positive, negative or zero. The number of positive numbers formed by this sum is equal to the number of negative numbers formed by this sum, because of symmetry. There is only one way to achieve a sum of zero, if all $a_i=0$. The total number of ways to pick $a_i$ from $i=0, 1, 2, 3, ... 7$ is $3^8=6561$. $\frac{6561-1}{2}=3280$ gives the number of possible negative integers. The question asks for the number of non-negative integers, so subtracting from the total gives $6561-3280=\boxed{\textbf{(D) } 3281}$. (RegularHexagon, KLBBC minor changes)

 



\textbf{Solution 3 (Answer choices)}

Note that the number of total possibilities (ignoring the conditions set by the problem) is $3^8=6561$. So, E is clearly unrealistic. 

 Note that if $a_7$ is 1, then it's impossible for 
 \[a_7\cdot3^7+a_6\cdot3^6+a_5\cdot3^5+a_4\cdot3^4+a_3\cdot3^3+a_2\cdot3^2+a_1\cdot3^1+a_0\cdot3^0,\] to be negative. Therefore, if $a_7$ is 1, there are $3^7=2187$ possibilities. (We also must convince ourselves that these $2187$ different sets of coefficients must necessarily yield $2187$ different integer results.)

 As A, B, and C are all less than 2187, the answer must be $\boxed{\textbf{(D) } 3281}$




\textbf{Solution 4}

Note that we can do some simple casework:If $a_7=1$, then we can choose anything for the other 7 variables, so this give us $3^7$.If $a_7=0$ and $a_6=1$, then we can choose anything for the other 6 variables, giving us $3^6$.If $a_7=0$, $a_6=0$, and $a_5=1$, then we have $3^5$.Continuing in this vein, we have $3^7+3^6+\cdots+3^1+3^0$ ways to choose the variables' values, except we have to add 1because we haven't counted the case where all variables are 0. So our total sum is $\boxed{\textbf{(D) } 3281}$.Note that we have counted all possibilities, because the largest positive positive power of 3 must be greater than or equal to the largest negative positive power of 3, for the number to be nonnegative.

 



\textbf{Solution 5 (Casework and Recursion)}

This solution has a similar idea to that of above.We will start off with casework on $a_{7}$:If $a_7 = -1$, we can show that there is no way to assign the values of $-1, 0, 1$ to the other $a_{i}$'s for $0 \leq i \leq 6$. This is because even if we make all of the other $a_{i}$ equal to $1$, we will have that our number is $(-1)(3^{7}) + (3^{6}+3^{5}+...+3^{0}) = -3^{7} + \frac{3^{7}-1}{3-1}$. However, note that since $\frac{3^{7}-1}{3-1} < 3^{7}$, $-3^{7} + \frac{3^{7}-1}{3-1} < 0$.If $a_{7}=0$, we can ignore $a_{7}$ and pretend like it never even existed. This then simplifies the problem: the question remains the same, except that instead of having $a_{0}...a_{7}$, we only have $a_{0}...a_{6}$. Now, we introduce some notation: let the number of nonnegative integers that can be written in the form of $a_{n} \cdot 3^{n} + ... + a_{0} \cdot 3^{0}$ be denoted as $S_{n}$. Then, the problem is asking us for $S_{7}$, and the total number of nonnegative integers we get from this case is $S_{6}$. Let's keep this in mind and revisit this idea later.If $a_{7}=1$, we can show that no matter what values we assign to the rest of the $a_{i}$, we will always get a nonnegative integer (we can make this even stricter and say positive integer). This is because of something we figured out in the first case: $\frac{3^{7}-1}{3-1} < 3^{7}$. So, even if we let $a_{i}=-1$ for $0 \leq i \leq 6$, we will still have $3^{7} - (3^{6}+...+3^{0}) > 0$. So, for this case, we have $3^7$ possible nonnegative integers.Totalling up our values from the three cases, we see that $S_{7} = 0 + S_{6} + 3^{7}$. In fact, using the reasoning from the three cases above, we can deduce that, for all positive integers $n$, $S_{n} = 3^{n} + S_{n-1}$. We can now start building up our values for $S$: $S_{0} = 2$ ($a_{0} = 0, 1$), $S_{1} = 3^{1}+S_{0} = 5, S_{2} = 14, ..., S_{6}=1094$. So, we have $S_{7} = 2187 + 1094 = \boxed{3281}$, which is answer choice $\boxed{\text{D}}$.~advanture

 



\textbf{Solution 6}

The key is to realize that this question is basically taking place in $a\in\{0,1,2\}$ if each value of $a$ was increased by $1$, essentially making it into base $3$. Then the range would be from $0\cdot3^7+$ $0\cdot3^6+$ $0\cdot3^5+$ $0\cdot3^4+$ $0\cdot3^3+$ $0\cdot3^2+$ $0\cdot3^1+$ $0\cdot3^0=$ $0$ to $2\cdot3^7+$ $2\cdot3^6+$ $2\cdot3^5+$ $2\cdot3^4+$ $2\cdot3^3+$ $2\cdot3^2+$ $2\cdot3^1+$ $2\cdot3^0=$ $3^8-1=$ $6561-1=$ $6560$, yielding $6561$ different values. Since the distribution for all $a_i\in \{-1,0,1\}$ the question originally gave is symmetrical, we retain the $3280$ positive integers and one $0$ but discard the $3280$ negative integers. Thus, we are left with the answer, $\boxed{\textbf{(D)} 3281}\qquad$.  --anna0kear

 



\textbf{Solution 7}

First, set $a_i=0$ for all $i\geq1$. The range would be the integers for which $[-1,1]$. If $a_i=0$ for all $i\geq2$, our set expands to include all integers such that $-4\leq\mathbb{Z}\leq4$. Similarly, when $i\geq3$ we get $-13\leq\mathbb{Z}\leq13$, and when $i\geq4$ the range is $-40\leq\mathbb{Z}\leq40$. The pattern continues until we reach $i=7$, where $-3280\leq\mathbb{Z}\leq3280$. Because we are only looking for positive integers, we filter out all $\mathbb{Z}<0$, leaving us with all integers between $0\leq\mathbb{Z}\leq3280$, inclusive. The answer becomes $\boxed{\textbf{(D) } 3281}$.  --anna0kear

 



\textbf{Solution 8}

To get the number of integers, we can get the highest positive integer that can be represented using 
 \[a_7\cdot3^7+a_6\cdot3^6+a_5\cdot3^5+a_4\cdot3^4+a_3\cdot3^3+a_2\cdot3^2+a_1\cdot3^1+a_0\cdot3^0,\]where $a_i\in \{-1,0,1\}$ for $0\le i \le 7$.

 Note that the least nonnegative integer that can be represented is $0$, when all $a_i=0$. The highest number will be the number when all $a_i=1$. That will be 
 \[3^7+3^6+3^5+3^4+3^3+3^2+3^1+3^0=\frac{3^8-1}{3-1}\] 
 \[=3280\]

 Therefore, there are $3280$ positive integers and $(3280+1)$ nonnegative integers (while including $0$) that can be represented. Our answer is $\boxed{\textbf{(D) } 3281}$




\textbf{Solution 9}

Notice that there are $3^8$ options for $a_7, a_6, \cdots a_0$ since each $a_i$ can take on the value $-1$, $0$, or $1$. Now we want to find how many of them are positive and then we can add one in the end to account for $0$ (they are asking for non-negative). 

 By symmetry (look out for these on the contest), we see that exactly half of them are positive. So $\lfloor{\frac{3^8}{2}}\rfloor = 3280.$ Now we will add $1$ because of the $0$ to account for the non-negative solutions. 

 So our final answer is $3280 + 1 = 3281$ which is $\boxed{\textbf{(D) } 3281}$.

 



\textbf{Solution 10}

Obviously, there are $3^8 = 6561$ possible, and one of them is 0, so other $6560$ are either positive or negative. By the symmetry, we can get the answer is $6560/2 + 1$ which is $\boxed{\textbf{(D) } 3281}$.

 



\textbf{Solution 11 (2 seconds)}

The solution has to obviously be smaller than $3^{7-0+1}=3^8$ but larger than $3^{(7-0+1)-1} = 3^7.$ The only remaining answer is $\boxed{\textbf{(D) } 3281}.$~mathboy282

 



\textbf{Solution 12 (Sets)}

Define Set $S=\{3^0, 3^1, 3^2, 3^3, 3^4, 3^5, 3^6, 3^7\}$.

 Suppose $T$ is a subset of $S$ with $k$ values. Then we can express $\sum_{m=0}^{7} a_m3^m$ as
 \[\sum_{d \in T} \pm d\]

 Even though we can arbitrarily assign $a_d$ for all $d \in T$ as either $-1$ or $1$, because $\sum_{d \in T} \pm d$ needs to be nonnegative, for the largest element $n$ of $T$ $a_n$ must be equal to $1$. This means that for any subset $T$ with $k$ values there are $2^{k-1}$ ways to assign $a_d$ for all $d \in T$.

 Moreover, there are $\binom{8}{k}$ subsets $T$ of $S$ with $k$ values. So our total number of positive numbers that could be expressed as $\sum_{m=0}^{7} a_m3^m$ is
 \[\sum_{k=1}^{8} \binom{8}{k}2^{k-1}=3280.\]

 But wait---the problem is asking for the number of nonnegative numbers, not just positive. So we need to find the cases where $\sum_{m=0}^{7} a_m3^m=0$. Fortunately, this only happens when $T=\emptyset$, so there is only $1$ case where this could happen. This could be verified with the fact that $\sum_{m=0}^{7} a_m3^m=0$ only occurs when $\{a_0, a_1, a_2, a_3, a_4, a_5, a_6, a_7\}=0$.Therefore, the answer is $3280+1=\boxed{\textbf{(D) }3281}$.

 \textbf{Soution 13 (Answer Choices)} Let's look at the number we form. If we can make the number $x$, the we can also make $-x$. This would be an even amount of total cases, except it causes $0$ and $-0$ to be marked as different numbers. We must then subtract one from our final value. Adding one to all our options, we are left with $\text{(B)}$ and $\text{(D)}$. It's not $\text{(B)}$ because it's $3^6$, and there are 8 sections in the equation, so therefore our answer must be $\boxed{\textbf{(D) }3281}$.

 


	\item (\textit{AMC 12 12A 2018}) The solution to the equation $\log_{3x} 4 = \log_{2x} 8$, where $x$ is a positive real number other than $\frac{1}{3}$ or $\frac{1}{2}$, can be written as $\frac {p}{q}$ where $p$ and $q$ are relatively prime positive integers. What is $p + q$?


 $\textbf{(A) } 5   \qquad     \textbf{(B) } 13   \qquad     \textbf{(C) } 17   \qquad    \textbf{(D) } 31 \qquad   \textbf{(E) } 35$

 \textbf{Answer: }D



\textbf{Solution 1}

We apply the Change of Base Formula, then rearrange:
 \begin{align*} \frac{\log_2{4}}{\log_2{(3x)}}&=\frac{\log_2{8}}{\log_2{(2x)}} \\ \frac{2}{\log_2{(3x)}}&=\frac{3}{\log_2{(2x)}} \\ 3\log_2{(3x)}&=2\log_2{(2x)}. \\ \end{align*}
By the logarithmic identity $n\log_b{a}=\log_b{\left(a^n\right)},$ it follows that
 \begin{align*} \log_2{\left[(3x)^3\right]}&=\log_2{\left[(2x)^2\right]} \\ (3x)^3&=(2x)^2 \\ 27x^3&=4x^2 \\ x&=\frac{4}{27}, \end{align*}
from which the answer is $4+27=\boxed{\textbf{(D) } 31}.$




\textbf{Solution 2}

We will apply the following logarithmic identity:
 \[\log_{p^n}{\left(q^n\right)}=\log_{p}{q},\]which can be proven by the Change of Base Formula: 
 \[\log_{p^n}{\left(q^n\right)}=\frac{\log_{p}{\left(q^n\right)}}{\log_{p}{\left(p^n\right)}}=\frac{n\log_{p}{q}}{n}=\log_{p}{q}.\]We rewrite the original equation as $\log_{(3x)^3} 64 = \log_{(2x)^2} 64,$ from which
 \begin{align*} (3x)^3&=(2x)^2 \\ 27x^3&=4x^2 \\ x&=\frac{4}{27}. \end{align*}
Therefore, the answer is $4+27=\boxed{\textbf{(D) } 31}.$




\textbf{Solution 3}

By the logarithmic identity $n\log_b{a}=\log_b{\left(a^n\right)},$ the original equation becomes 
 \[2\log_{3x} 2 = 3\log_{2x} 2.\]By the logarithmic identity $\log_b{a}\cdot\log_a{b}=1,$ we multiply both sides by $\log_2{(2x)},$ then apply the Change of Base Formula to the left side:
 \begin{align*} 2\left[\log_{3x}2\right]\left[\log_2{(2x)}\right] &= 3 \\ 2\left[\frac{\log_2 2}{\log_2{(3x)}}\right]\left[\frac{\log_2{(2x)}}{\log_2 2}\right] &= 3 \\ 2\left[\frac{\log_2{(2x)}}{\log_2{(3x)}}\right] &=3 \\ 2\left[\log_{3x}{(2x)}\right] &= 3 \\ \log_{3x}{\left[(2x)^2\right]} &= 3 \\ (3x)^3&=(2x)^2 \\ 27x^3&=4x^2 \\ x&=\frac{4}{27}. \end{align*}
Therefore, the answer is $4+27=\boxed{\textbf{(D) } 31}.$




\textbf{Solution 4}

We can convert both $4$ and $8$ into $2^2$ and $2^3,$ respectively: 
 \[2\log_{3x} (2) = 3\log_{2x} (2).\]Converting the bases of the right side, we get
 \begin{align*} \log_{2x} 2 &= \frac{\ln 2}{\ln (2x)} \\ \frac{2}{3}\cdot\log_{3x} (2) &= \frac{\ln 2}{\ln (2x)} \\ 2^\frac{2}{3} &= (3x)^\frac{\ln 2}{\ln (2x)} \\ \frac{2}{3} \cdot \ln 2 &= \frac{\ln 2}{\ln (2x)} \cdot \ln (3x). \end{align*}
Dividing both sides by $\ln 2,$ we get $\frac{2}{3} = \frac{\ln (3x)}{\ln (2x)},$ from which 
 \[2\ln (2x) = 3\ln (3x).\]Expanding this equation gives
 \begin{align*} 2\ln 2 + 2\ln (x) &= 3\ln 3 + 3\ln (x) \\ \ln (x) &= 2\ln 2 - 3\ln 3. \end{align*}
Thus, we have 
 \[x = e^{2\ln 2 - 3\ln 3} = \frac{e^{2\ln 2}}{e^{3\ln 3}} = \frac{2^2}{3^3} = \frac{4}{27},\]from which the answer is $4+27=\boxed{\textbf{(D) } 31}.$




\textbf{Solution 5 (Exponential Form)}

Let $y=\log_{3x} 4 = \log_{2x} 8.$ We convert the equations with $y$ to the exponential form:
 \begin{align*} (3x)^y&=4, \\ (2x)^y&=8. \end{align*}
Cubing the first equation and squaring the second equation, we have
 \begin{align*} (3x)^{3y}&=64, \\ (2x)^{2y}&=64. \end{align*}
Applying the Transitive Property, we get 
 \begin{align*} (3x)^{3y}&=(2x)^{2y} \\ (3x)^3&=(2x)^2 \\ 27x^3&=4x^2 \\ x&=\frac{4}{27}, \end{align*}
from which the answer is $4+27=\boxed{\textbf{(D) } 31}.$



	\item (\textit{AMC 12 12A 2018}) A scanning code consists of a $7 \times 7$ grid of squares, with some of its squares colored black and the rest colored white. There must be at least one square of each color in this grid of $49$ squares. A scanning code is called $\textit{symmetric}$ if its look does not change when the entire square is rotated by a multiple of $90 ^{\circ}$ counterclockwise around its center, nor when it is reflected across a line joining opposite corners or a line joining midpoints of opposite sides. What is the total number of possible symmetric scanning codes?


 $\textbf{(A)} \text{ 510} \qquad \textbf{(B)} \text{ 1022} \qquad \textbf{(C)} \text{ 8190} \qquad \textbf{(D)} \text{ 8192} \qquad \textbf{(E)} \text{ 65,534}$

 \textbf{Answer: }B



\textbf{Solution 1}

Draw a $7 \times 7$ square.

 
 \[\begin{tabular}{|c|c|c|c|c|c|c|} \hline K & J & H & G & H & J & K \\ \hline J & F & E & D & E & F & J \\ \hline H & E & C & B & C & E & H \\ \hline G & D & B & A & B & D & G \\ \hline H & E & C & B & C & E & H \\ \hline J & F & E & D & E & F & J \\ \hline K & J & H & G & H & J & K \\ \hline \end{tabular}\]

 Start from the center and label all protruding cells symmetrically. (Note that "I" is left out of this labelling, so there are only 10 labels, not 11, as ending in K would suggest.)

 More specifically, since there are $4$ given lines of symmetry ($2$ diagonals, $1$ vertical, $1$ horizontal) and they split the plot into $8$ equivalent sections, we can take just one-eighth and study it in particular. Each of these sections has $10$ distinct sub-squares, whether partially or in full. So since each can be colored either white or black, we choose $2^{10}=1024$ but then subtract the $2$ cases where all are white or all are black. That leaves us with $\fbox{\textbf{(B)} \text{ 1022}}$.

 There are only ten squares we get to actually choose, and two independent choices for each, for a total of $2^{10} = 1024$ codes. Two codes must be subtracted (due to the rule that there must be at least one square of each color) for an answer of $\fbox{\textbf{(B) }1022}$.

 
Note that this problem is very similar to the 1996 AIME Problem 7 \href{https://artofproblemsolving.com/wiki/index.php/1996\_AIME\_Problems/Problem\_7}{https://artofproblemsolving.com/wiki/index.php/1996\_AIME\_Problems/Problem\_7}.

 



\textbf{Solution 2 (Very Intuitive and Fast)}


                    \begin{figure}[H]
                        \centering
                        \includegraphics[width=0.5\textwidth]{Solution/AMC_12_12A_2018_15_Solution_Figure_1.png}
                    \end{figure}
                    

 Imagine folding the scanning code along its lines of symmetry. There will be $10$ regions which you have control over coloring. Since we must subtract off $2$ cases for the all-black and all-white cases, the answer is $2^{10}-2=\boxed{\textbf{(B) } 1022.}$




\textbf{Solution 3}


 \[\begin{tabular}{|c|c|c|c|c|c|c|} \hline T & T & T & X & T & T & T \\ \hline T & T & T & Y& T & T & T \\ \hline T & T & T & Z & T & T & T \\ \hline X & Y & Z & W & Z & Y & X \\ \hline T & T & T & Z & T & T & T \\ \hline T & T & T & Y & T & T & T \\ \hline T & T & T & X & T & T & T \\ \hline \end{tabular}\]

 There are $3 \times 3$ squares in the corners of this $7 \times 7$ square, and there is a horizontal and vertical stripe through the middle. Because we need to have symmetry when the diagonals and midpoints of the large square is connected, we can create a table like this: (Different letters represent different color choices between black and white)

 
 \[\begin{tabular}{|c|c|c|} \hline A & B & C \\ \hline B & D & E \\ \hline C & E & F \\ \hline \end{tabular}\](Note that coloring one $3 \times 3$ square will also determine the colorings of the other $3$ because the large $7 \times 7$ square must look the same when it is rotated by $90^\circ$)

 There are $6$ different letters and $2$ choices of color (black and white) for each letter, so there are $2^6=64$ colorings of a proper $3 \times 3$ square. Now, all that's left are the horizontal and vertical stripes. Using similar logic, we can see that there are $4$ different letters, so there are $2^4=16$ different colorings. Multiplying them together gives $16 \times 64 = 1024$. Going back to the question, we see that "there must be at least one square of each color in this grid of $49$ squares." We must then eliminate $2$ options: an all-black grid and an all-white grid. $1024-2= \boxed{\textbf{(B)} 1022} \\ \phantom{}$ -hansenhe

 



\textbf{Solution 4 (Similar to Solution 2)}

A symmetric code following the second and third conditions (diagonal and straight lines of symmetry) will automatically fulfill the first one. We can work with the grid by splitting it into different categories: the center (C), the diagonals across the corners (D), the straight lines across the midpoints (S), and the remaining cubes (R) as shown in the diagram below.

 
 \[\begin{tabular}{|c|c|c|c|c|c|c|} \hline D & R & R & S & R & R & D \\ \hline R & D & R & S & R & D & R \\ \hline R & R & D & S & D & R & R \\ \hline S & S & S & C & S & S & S \\ \hline R & R & D & S & D & R & R \\ \hline R & D & R & S & R & D & R \\ \hline D & R & R & S & R & R & D \\ \hline \end{tabular}\]

 If we account for symmetry, we will only need to pay attention to a small section:

 
 \[\begin{tabular}{|c|c|c|c|} \hline D &   &   &   \\ \hline R & D &   &   \\ \hline R & R & D &   \\ \hline S & S & S & C \\ \hline \end{tabular}\] 

 Note: this shows part of the full grid

 There are $2^3$ states of D, R, and S, and there are 2 states of C so $2^3 \cdot 2^3 \cdot 2^3 \cdot 2$ and subtracting 2 cases for all white and all black will bring $2^{10} - 2 = \boxed{\textbf{(B)} 1022}$ as our answer.

 


	\item (\textit{AMC 12 12A 2018}) Which of the following describes the set of values of $a$ for which the curves $x^2+y^2=a^2$ and $y=x^2-a$ in the real $xy$-plane intersect at exactly $3$ points?


 $\textbf{(A) }a=\frac14 \qquad \textbf{(B) }\frac14 < a < \frac12 \qquad \textbf{(C) }a>\frac14 \qquad \textbf{(D) }a=\frac12 \qquad \textbf{(E) }a>\frac12 \qquad$

 \textbf{Answer: }E



\textbf{Solution 1 (Algebra)}

Substituting $y=x^2-a$ into $x^2+y^2=a^2$, we get
 \[x^2+(x^2-a)^2=a^2 \implies x^2+x^4-2ax^2=0 \implies x^2(x^2-(2a-1))=0\]Since this is a quartic, there are $4$ total roots (counting multiplicity). We see that $x=0$ always has at least one intersection at $(0,-a)$ (and is in fact a double root). 

 The other two intersection points have $x$ coordinates $\pm\sqrt{2a-1}$. We must have $2a-1> 0;$ otherwise we are in the case where the parabola lies entirely below the circle (tangent at the point $(0,-a)$). This only results in a single intersection point in the real coordinate plane. Thus, we see that $\boxed{\textbf{(E) }a>\frac12}$.

 (projecteulerlover)

 



\textbf{Solution 2 (Algebra)}

Substituting $y = x^2 - a$ gives $x^2 + (x^2 - a)^2 = a^2$, which simplifies to $x^2 + x^4 - 2x^2a + a^2 = a^2$. This further simplifies to $x^2(1 + x^2 - 2a) = 0$. Thus, either $x^2 = 0$, or $x^2 - 2a + 1 = 0$. 

 Since we care about $a$, we consider the second case. We solve in terms of $a$, giving $a = \frac{x^2}{2} + \frac{1}{2}$. We see that in order to find the range in which $a$ lies, we must find the vertex of this equation, which turns out to be $\left(0, \frac{1}{2}\right)$. Hence, we know that the minimum is $\frac{1}{2}$, which further implies that $\boxed{\textbf{(E) }a>\frac12}$.

 



\textbf{Solution 3 (Graphs)}


                    \begin{figure}[H]
                        \centering
                        \includegraphics[width=0.5\textwidth]{Solution/AMC_12_12A_2018_16_Solution_Figure_1.png}
                    \end{figure}
                    

 Looking at a graph, it is obvious that the two curves intersect at $(0, -a)$. We also see that if the parabola goes "in" the circle, then by going out of it (as it will), it will intersect five times. This is impossible. Thus, we only look for cases where the parabola becomes externally tangent to the circle. 

 We have $x^2 - a = -\sqrt{a^2 - x^2}$. Squaring both sides and solving yields $x^4 - (2a - 1)x^2 = 0$. Since $x = 0$ is already accounted for, we only need to find one solution for $x^2 = 2a - 1$, where the right hand side portion is obviously increasing. Since $a = \frac{1}{2}$ begets $x = 0$ (an overcount), we have $\boxed{\textbf{(E) }a>\frac12}$ as the right answer.

 Solution by JohnHankock

 



\textbf{Solution 4 (Graphs)}

This describes a unit parabola, with a circle centered on the axis of symmetry and tangent to the vertex. As the curvature of the unit parabola at the vertex is $2$, the radius of the circle that matches it has a radius of $\frac{1}{2}$. This circle is tangent to an infinitesimally close pair of points, one on each side. Therefore, it is tangent to only $1$ point. When a larger circle is used, it is tangent to $3$ points because the points on either side are now separated from the vertex. Therefore, $\boxed{\textbf{(E) }a>\frac12}$ is correct.

 



\textbf{Solution 5 (Graphs)}

Now, let's graph these two equations. We want the blue parabola to be inside this red circle.
                    \begin{figure}[H]
                        \centering
                        \includegraphics[width=0.5\textwidth]{Solution/AMC_12_12A_2018_16_Solution_Figure_2.png}
                    \end{figure}
                    Then we substitute $y$ into the first equation to get $x^2+(x^2-a)^2=a^2$. Expanding, we get $x^4-2ax^2+x^2=0$. Factoring out the $x^2$, we get $x^2(x^2-2a+1)=0$. Then we find that $x=0$ or $x=\pm\sqrt{2a-1}$. Therefore, $2a-1>0$, which means $\boxed{\textbf{(E) }a>\frac12}$.

 - kante314 -~ minor edit by junokim1011 and RJ5303707

 



\textbf{Solution 6 (Calculus)}

In order to solve for the values of $a$, we need to just count multiplicities of the roots when the equations are set equal to each other: in other words, take the derivative. We know that $\sqrt{a^2 - x^2} = x^2 - a$. Now, we take square of both sides, and rearrange to obtain $x^4 - (2a - 1)x^2 = 0$. Now, we may take the second derivative of the equation to obtain $6x^2 - (2a - 1) = 0$. Now, we must take discriminant. Since we need the roots of that equation to be real and not repetitive (otherwise they would not intersect each other at three points), the discriminant must be greater than zero. Thus,

 
 \[b^2 - 4ac > 0 \rightarrow 0 - 4(6)(-(2a - 1)) > 0 \rightarrow a > \frac{1}{2}\]The answer is $\boxed{\textbf{(E) }a>\frac12}$ and we are done. 

 (Edited by OlutosinNGA)

 



\textbf{Solution 7 (Calculus)}

Notice, the equations are of that of a circle of radius a centered at the origin and a parabola translated down by a units. They always intersect at the point $(0, a)$, and they have symmetry across the $y$-axis, thus, for them to intersect at exactly $3$ points, it suffices to find the $y$ solution. 

 First, rewrite the second equation to $y=x^2-a\implies x^2=y+a$And substitute into the first equation: $y+a+y^2=a^2$ Since we're only interested in seeing the interval in which a can exist, we find the discriminant: $1-4a+4a^2$. This value must not be less than $0$ (It is the square root part of the quadratic formula). To find when it is $0$, we find the roots: 
 \[4a^2-4a+1=0 \implies a=\frac{4\pm\sqrt{16-16}}{8}=\frac{1}{2}\]Since $\lim_{a\to \infty}(4a^2-4a+1)=\infty$, our range is $\boxed{\textbf{(E) }a>\frac12}$.

 Solution by ktong

 



\textbf{Solution 8 (Observations)}

We can see that if $a = 1$, we know that the points where the two curves intersect are $(0, -1), (1, 0)$ and $(-1, 0)$ .Because there are only $3$ intersections and $a > 1/2$, as well as $a > 1/4$ we know that either $\textbf{(D)}$ or $\textbf{(E)}$ is the correct answer. Then we can test a number from $(1/2, 1/4)$ to eliminate the remaining answer. So  $\boxed{\textbf{(E) }a>\frac12}$ is the correct answer. 

 



\textbf{Solution 9 (Observations)}

Simply plug in $a = \frac{1}{2}, \frac{1}{4}, 1$ and solve the systems. (This shouldn't take too long.) And then realize that only $a=1$ yields three real solutions for $x$, so we are done and the answer is $\boxed{\textbf{(E) }a>\frac12}$.

 



\textbf{Solution 10 (Observations)}

An ideal solution come to mind is where they intersect at the $x$-axis at the same time, which is $(a, 0)$ and $(-a, 0).$ Take the root of our $y=x^2-a$ we get $x=\sqrt{a}$, set them equal we get $a=\sqrt{a}.$ The only answer is $1$ so it only left us with the answer choice $\boxed{\textbf{(E) }a>\frac12}$.

 



\textbf{Solution 11 (Curvature, Calculus)}

From observing the graph of the two functions, we notice that there will be exactly 3 solutions when the curvature of $y = x^2-a$ is greater than the curvature of the circle $x^2+y^2 = a^2$ near $x=0$ (this means the graph of $y=x^2-a$ will reside within the circle). Formally, curvature is the magnitude of the rate of change of the unit tangent in respect to arc length. Let there be function $r(x)$ be a vector-valued function. The unit tangent has formula $T(x)=\frac{r'(x)}{|r'(x)|}=\frac{dr}{ds}$, and arc length has formula $s(t)=\int_b^{t}\sqrt{\left( 1+ \left( \frac{dr}{dx} \right) ^2 \right) }dx$. We have the relation $\frac{ds}{dx}=\sqrt{\left( 1+ \left( \frac{dr}{dx} \right) ^2 \right) }=|r'(x)|$. Our goal, curvature, is $\kappa= \left| \frac{d^2r}{ds^2} \right| =\frac{dT}{ds}|=\frac1{|\frac{ds}{dt}|}|\frac{dT}{dt}|$. After deriving, we end up with the formula:
 \[\kappa = \frac{y^{\prime\prime}}{(1+(y^{\prime})^2)^\frac{3}{2}}\]For a circle, curvature is just the inverse of the radius or $1/a$ in our case. For the parabola, we use the equation above:
 \[\kappa = \frac{y^{\prime\prime}}{(1+(y^{\prime})^2)} = \frac{2}{(1+4x^2)^{1.5}}\]where $\kappa \approx 2$ when $x \approx 0$. Since curvature of the parabola should be greater than that of the circle, $2>\frac{1}{a}$ or answer choice $\boxed{\textbf{(E)}a>\frac12}$.

 \textbf{Remark} The graph of $x^2+y^2=a^2$ is a circle with the center $(0,0)$ and the radius $|a|.$

 The graph of $y=x^2-a$ is an upward-opening parabola with the vertex $(0,-a).$

 The circle and the parabola are shown here in Desmos: \href{https://www.desmos.com/calculator/rfeknuhjum}{https://www.desmos.com/calculator/rfeknuhjum}

 Move the slider around for $\frac12\leq a\leq5$ to observe that they intersect at exactly $3$ points: $(0,-a)$ and $\left(\pm\sqrt{2a-1},a-1\right),$ provided that $a>\frac12.$




\textbf{Solution 12 (cheap)}

We observe that there is more than one value for $a$, so we can eliminate $(A)$ and $(D)$. We can also observe that as we increase $a$, the parabola will always intersect the circle at 3 points, which eliminates $(B)$. Finally, we notice that if the $a=1/4$ case is not a perfect square, then $(C)$ is not the correct answer.

 


 Substituting $1/4$ for $a$, we get:

 


 $x^2+y^2=(1/4)^2$
$x^2-1/4=y$ --> $x^2=y+1/4$
$y^2+y+1/4=(1/4)^2$ --> $y^2+y+3/16$

 


 From this, we can see that the $a=1/4$ case is incorrect, so our final answer is $\boxed{\textbf{(E)}a>\frac12}$.

 - WB0s0n

 


	\item (\textit{AMC 12 12A 2018}) Farmer Pythagoras has a field in the shape of a right triangle. The right triangle's legs have lengths $3$ and $4$ units. In the corner where those sides meet at a right angle, he leaves a small unplanted square $S$ so that from the air it looks like the right angle symbol. The rest of the field is planted. The shortest distance from $S$ to the hypotenuse is $2$ units. What fraction of the field is planted?


 
                    \begin{figure}[H]
                        \centering
                        \includegraphics[width=0.5\textwidth]{Problem/AMC_12_12A_2018_17_Problem_Figure_1.png}
                    \end{figure}
                    
$\textbf{(A) }   \frac{25}{27}   \qquad        \textbf{(B) }   \frac{26}{27}   \qquad    \textbf{(C) }   \frac{73}{75}   \qquad   \textbf{(D) } \frac{145}{147} \qquad  \textbf{(E) }   \frac{74}{75}$

 \textbf{Answer: }D



\textbf{Solution 1 (Area Addition)}

Note that the hypotenuse of the field is $5,$ and the area of the field is $6.$ Let $x$ be the side-length of square $S.$

 We partition the field into a red triangle, a yellow triangle, and a green triangle, as shown below:
                    \begin{figure}[H]
                        \centering
                        \includegraphics[width=0.5\textwidth]{Solution/AMC_12_12A_2018_17_Solution_Figure_1.png}
                    \end{figure}
                    Let the brackets denote areas. By area addition, we set up an equation for $x:$
 \begin{align*} [\text{Red Triangle}]+[\text{Yellow Triangle}]+[\text{Green Triangle}]&=[\text{Field}] \\ \frac{3x}{2}+\frac{4x}{2}+\frac{5\cdot2}{2}&=6, \end{align*}
from which $x=\frac27.$ Therefore, the answer is 
 \[\frac{[\text{Field}]-[S]}{[\text{Field}]}=\frac{6-x^2}{6}=\boxed{\textbf{(D) } \frac{145}{147}}.\]~MRENTHUSIASM

 



\textbf{Solution 2 (Area Addition)}

Let the square have side length $x$. Connect the upper-right vertex of square $S$ with the two vertices of the triangle's hypotenuse. This divides the triangle in several regions whose areas must add up to the area of the whole triangle, which is $6$.
                    \begin{figure}[H]
                        \centering
                        \includegraphics[width=0.5\textwidth]{Solution/AMC_12_12A_2018_17_Solution_Figure_2.png}
                    \end{figure}
                    Square $S$ has area $x^2$, and the two thin triangle regions have area $\dfrac{x(3-x)}{2}$ and $\dfrac{x(4-x)}{2}$. The final triangular region with the hypotenuse as its base and height $2$ has area $5$. Thus, we have 
 \[x^2+\dfrac{x(3-x)}{2}+\dfrac{x(4-x)}{2}+5=6.\]Solving gives $x=\dfrac{2}{7}$. The area of $S$ is $\dfrac{4}{49}$ and the desired ratio is $\dfrac{6-\frac{4}{49}}{6}=\boxed{\textbf{(D) } \frac{145}{147}}$.

 Alternatively, once you get $x=\frac{2}{7}$, you can avoid computation by noticing that there is a denominator of $7$, so the answer must have a factor of $7$ in the denominator, which only $\frac{145}{147}$ does.

 



\textbf{Solution 3 (Similar Triangles)}

Let the square have side length $s$. If we were to extend the sides of the square further into the triangle until they intersect on point on the hypotenuse, we'd have a similar right triangle formed between the hypotenuse and the two new lines, and two smaller similar triangles that share a side of length $2$. Using the side-to-side ratios of these triangles, we can find that the length of the larger similar triangle is $\frac{5}{3}(2)=\frac{10}{3}$. Now, let's extend this larger similar right triangle to the left until it hits the side of length $3$. Now, the length is $\frac{10}{3}+s$, and using the ratios of the side lengths, the height is $\frac{3}{4}\left(\frac{10}{3}+s\right)=\frac{5}{2}+\frac{3s}{4}$. Looking at the diagram, if we add the height of this triangle to the side length of the square, we'd get $3$, so
 \begin{align*} \frac{5}{2}+\frac{3s}{4}+s&=3 \\  \frac{5}{2}+\frac{7s}{4}&=3 \\ \frac{7s}{4}&=\frac{1}{2} \\ s&=\frac{2}{7}. \end{align*}
So, the area of the square is $\left(\frac{2}{7}\right)^2=\frac{4}{49}$.

 Now comes the easy part--finding the ratio of the areas: $\frac{3\cdot 4 \cdot \frac{1}{2} -\frac{4}{49}}{3\cdot 4 \cdot \frac{1}{2}}=\frac{6-\frac{4}{49}}{6}=\frac{294-4}{294}=\frac{290}{294}=\boxed{\textbf{(D) } \frac{145}{147}}$.

 



\textbf{Solution 4 (Similar Triangles)}


                    \begin{figure}[H]
                        \centering
                        \includegraphics[width=0.5\textwidth]{Solution/AMC_12_12A_2018_17_Solution_Figure_3.png}
                    \end{figure}
                    On the diagram above, find two smaller triangles similar to the large one with side lengths $3$, $4$, and $5$; consequently, the segments with length $\frac{5}{2}$ and $\frac{10}{3}$.

 With $\ell$ being the side length of the square, we need to find an expression for $\ell$. Using the hypotenuse, we can see that $\frac{3}{2}+\frac{8}{3}+\frac{5}{4}\ell+\frac{5}{3}\ell=5$. Simplifying, $\frac{35}{12}\ell=\frac{5}{6}$, or $\ell=\frac27$.

 A different calculation would yield $\ell+\frac{3}{4}\ell+\frac{5}{2}=3$, so $\frac{7}{4}\ell=\frac{1}{2}$. In other words, $\ell=\frac{2}{7}$, while to check, $\ell+\frac{4}{3}\ell+\frac{10}{3}=4$. As such, $\frac{7}{3}\ell=\frac{2}{3}$, and $\ell=\frac{2}{7}$.

 Finally, we get $A(\Square S)=\ell^2=\frac{4}{49}$, to finish. As a proportion of the triangle with area $6$, the answer would be $1-\frac{4}{49\cdot6}=1-\frac{2}{147}=\frac{145}{147}$, so $\boxed{\textbf{(D) } \frac{145}{147}}$ is correct.

 



\textbf{Solution 5 (Similar Triangles)}

Let the side length of the square be $x$. First off, let us make a similar triangle with the segment of length $2$ and the top-right corner of $S$. Therefore, the longest side of the smaller triangle must be $2 \cdot \frac54 = \frac52$. We then do operations with that side in terms of $x$. We subtract $x$ from the bottom, and $\frac{3x}{4}$ from the top. That gives us the equation of $3-\frac{7x}{4} = \frac{5}{2}$. Solving, 
 \[12-7x = 10 \implies x = \frac{2}{7}.\]Thus, $x^2 = \frac{4}{49}$, so the fraction of the triangle (area $6$) covered by the square is $\frac{2}{147}$. The answer is then $\boxed{\textbf{(D) } \frac{145}{147}}$.

 



\textbf{Solution 6 (Similar Triangles)}


                    \begin{figure}[H]
                        \centering
                        \includegraphics[width=0.5\textwidth]{Solution/AMC_12_12A_2018_17_Solution_Figure_4.png}
                    \end{figure}
                    Let $AD=x$. Note that $\triangle DEF$ is a $3{-}4{-}5$ triangle, so $EF=\frac{5}{4}x$ and $FD=\frac{3}{4}x$. $BF=BD+FD=4-x+\frac{3}{4}x=4-\frac{1}{4}x$. We know that $GE$ is $2$ from the problem so $GF=2+\frac{5}{4}x$. $\triangle FGB$ is also a $3{-}4{-}5$ triangle with $GF:BF=3:5$. We now have $3\left(4-\frac{1}{4}x\right)=5\left(2+\frac{5}{4}x\right)$. Solving this equation, we get that $x=\frac{2}{7}$ so the area of $S$ is $\frac{4}{49}$. The area of the triangle is $\frac{3\cdot 4}{2}=6$ so the fraction of field that is unplanted is $\frac{\frac{4}{49}}{6}=\frac{2}{147}$. Thus, the fraction of the field that is planted is $1-\frac{2}{147}=\boxed{\textbf{(D) } \frac{145}{147}}$.

 



\textbf{Solution 7 (Coordinate Geometry)}

We use coordinate geometry. Let the right angle be at $(0,0)$ and the hypotenuse be the line $3x+4y = 12$ for $0\le x\le 3$. Denote the position of $S$ as $(s,s)$, and by the point to line distance formula, we know that
 \begin{align*} \frac{|3s+4s-12|}{5} &= 2 \\ |7s-12| &= 10 \end{align*}
Solving this, we get $s=\frac{22}{7}, \frac{2}{7}$. Obviously $s<\frac{22}{7}$, so $s = \frac{2}{7}$, and from here, the rest of the solution follows to get $\boxed{\textbf{(D) } \frac{145}{147}}$.

 



\textbf{Solution 8 (Coordinate Geometry)}

Let the right angle be at $(0,0)$, the point $(x,x)$ be the far edge of the unplanted square and the hypotenuse be the line $y=-\frac{3}{4}x+3$. Since the line from $(x,x)$ to the hypotenuse is the shortest possible distance, we know this line, call it line $\l$, is perpendicular to the hypotenuse and therefore has a slope of $\frac{4}{3}$. 

 Since we know $m=\frac{4}{3}$ , we can see that the line rises by $\frac{8}{5}$ and moves to the right by $\frac{6}{5}$ to meet the hypotenuse. (Let $2 = 5x$ and the rise be $4x$ and the run be $3x$ and then solve.) Therefore, line $\l$ intersects the hypotenuse at the point $\left(x+\frac{6}{5}, x+\frac{8}{5}\right)$. Plugging into the equation for the hypotenuse we have $x=\frac{2}{7}$ , and after a bit of computation we get $\boxed{\textbf{(D) } \frac{145}{147}}$.

 



\textbf{Solution 9 (Pythagorean Theorem)}

Let the side length of the square be $a$, and the lengths that the line from $S$ hits the hypotenuse be $x$ and $5-x$. Also, connect the outermost vertex of $S$ to the vertices that $S$ isn't connected to. Note that the line that hits the hypotenuse must create a right angle, since it is the shortest possible distance. This creates two pairs of right triangles that share the same hypotenuse. This means that we can set up a system of equations using the Pythagorean Theorem: 
 \begin{align*} a^2+(4-a)^2&=(5-x)^2+2^2, \\ a^2+(3-a)^2&=x^2+2^2. \end{align*}
After *some* algebra, we obtain $a=\frac{2}{7}$, which gives the answer $\boxed{\textbf{(D) } \frac{145}{147}}$.

 



\textbf{Solution 10 (Proportions)}

We name <i>\textbf{small triangle}</i> the  triangle similar to given in which unplanted square $S$ is inscribed.The height of given triangle is 2.4 units so similarity coefficient is $\frac {2.4 - 2}{2.4} = \frac {1}{6}$ , the  area is $\frac {1}{36}$ of total area.

 The ratio of planted area in <i>\textbf{small triangle}</i> to the area of the square is $\frac {3}{8} + \frac {2}{3} = \frac {25}{24}.$

 The fraction of planted area in <i>\textbf{small triangle}</i> is $\frac {25}{25+24} = \frac {25}{49}.$

 Therefore, the fraction of the planted  field is $\frac {25}{49} \cdot  \frac {1}{36} + \frac {35}{36} = \boxed{\textbf{(D) } \frac{145}{147}}.$

 \textbf{vladimir.shelomovskii@gmail.com, vvsss}

 



\textbf{Solution 11 (Bash)}


                    \begin{figure}[H]
                        \centering
                        \includegraphics[width=0.5\textwidth]{Solution/AMC_12_12A_2018_17_Solution_Figure_5.png}
                    \end{figure}
                    

 Denote $A,B,C$ to be the three vertices of the triangular field. Also denote $A,M,D,N$ to be the vertices of the square $S$. Let $X$ be on $BC$ such that $AC\parallel DX$ and $Y$ be on $BC$ such that $AB\parallel DY$. Let $P$ and $Q$ be the foot of the altitudes from $X$ to $AC$ and from $Y$ to $AB$ respectively.

 Note that $\triangle ABC \sim \triangle DYX \sim \triangle PXC \sim \triangle QBY$. Thus, $PC = x \cdot \frac34$ and $QB = x \cdot \frac43$, making
 \begin{align*} DX &= 3-x-\dfrac{3}{4}x = 3-\dfrac{7}{4}x, \\ MQ &= 4-x-\dfrac{4}{3}x = 4-\dfrac{7}{3}x. \end{align*}
Also from the similarity ratio is the fact that $CX = \frac54 x$ and $BY = \frac53 x$, making
 \[XY = 5 - \dfrac{5}{4}x - \dfrac{5}{3}x = 5 - \dfrac{35}{12}x.\]Computing the area of $\triangle XDY$ in two ways gives an equation for $x$:
 \begin{align*} \left(3-\dfrac{7}{4}x\right)\left(4-\dfrac{7}{3}x\right) &= 2\cdot \left(5 - \dfrac{35}{12}x\right) \\ 10-\dfrac{35}{6}x &= \dfrac{49}{12}x^2 - 14x + 12 \\ \dfrac{49}{12}x^2 - \dfrac{49}{6}x + 2 &= 0 \\ 49x^2 - 98x + 24 &= 0 \\ x&=\dfrac{2}{7} \text{ or } \dfrac{12}{7}. \end{align*}
But $x=\dfrac{12}{7}$ is extraneous. Thus, the area of square $S = x^2 = \dfrac{4}{49}$, making the portion of the field that is planted being 
 \[1 - \dfrac{\frac{4}{49}}{6} = 1 - \dfrac{2}{147} = \boxed{\textbf{(D) } \frac{145}{147}}.\]

 -Solution by sml1809

 



\textbf{Solution 12}

Let the square have side length $x$. Note that when $x = \dfrac{12}{7}$, the square is inscribed and touching the hypotenuse. Denote $f(x)$ to be the minimum distance from the square to the hypotenuse. Notably, $f$ is linear with respect to $x$. $f(0) = \dfrac{12}{5}$, because it is simply the length of the hypotenuse's altitude. Similarly, $f(\dfrac{12}{7}) = 0$. We can find that $f(x) = \dfrac{12-7x}{5}$. Setting this equal to $2$, we get that $f(\dfrac{2}{7}) = 2$. Therefore, the side has side length $\dfrac27$, and has area $\dfrac{4}{49}$. So, the unshaded area is $1 - \dfrac{\frac{4}{49}}{6} = \boxed{\dfrac{145}{147}}$, or $(D)$. ~Puck\_0

 


	\item (\textit{AMC 12 12A 2018}) Triangle $ABC$ with $AB=50$ and $AC=10$ has area $120$. Let $D$ be the midpoint of $\overline{AB}$, and let $E$ be the midpoint of $\overline{AC}$. The angle bisector of $\angle BAC$ intersects $\overline{DE}$ and $\overline{BC}$ at $F$ and $G$, respectively. What is the area of quadrilateral $FDBG$?


 $\textbf{(A) }60 \qquad \textbf{(B) }65 \qquad \textbf{(C) }70 \qquad \textbf{(D) }75 \qquad \textbf{(E) }80 \qquad$

 \textbf{Answer: }D



\textbf{Solution 1}

Let $BC = a$, $BG = x$, $GC = y$, and the length of the perpendicular from $BC$ through $A$ be $h$. By angle bisector theorem, we have that 
 \[\frac{50}{x} = \frac{10}{y},\] where $y = -x+a$. Therefore substituting we have that $BG=\frac{5a}{6}$. By similar triangles, we have that $DF=\frac{5a}{12}$, and the height of this trapezoid is $\frac{h}{2}$. Then, we have that $\frac{ah}{2}=120$. We wish to compute $\frac{5a}{8}\cdot\frac{h}{2}$, and we have that it is $\boxed{\textbf{(D) }75}$ by substituting.

 



\textbf{Solution 2}

For this problem, we have $\triangle{ADE}\sim\triangle{ABC}$ because of SAS and $DE = \frac{BC}{2}$. Therefore, $\bigtriangleup ADE$ is a quarter of the area of $\bigtriangleup ABC$, which is $30$. Subsequently, we can compute the area of quadrilateral $BDEC$ to be $120 - 30 = 90$. Using the angle bisector theorem in the same fashion as the previous solution, we get that $\overline{BG}$ is $5$ times the length of $\overline{GC}$. We want the larger piece, as described by the problem. Because the heights are identical, one area is $5$ times the other, and $\frac{5}{6} \cdot 90 = \boxed{\textbf{(D) }75}$.

 



\textbf{Solution 3}

The ratio of the $\overline{BG}$ to $\overline{GC}$ is $5:1$ by the Angle Bisector Theorem, so area of $\bigtriangleup ABG$ to the area of $\bigtriangleup ACG$ is also $5:1$ (They have the same height). Therefore, the area of $\bigtriangleup ABG$ is $\frac{5}{5+1}\times120=100$. Since $\overline{DE}$ is the midsegment of $\bigtriangleup ABC$, so $\overline{DF}$ is the midsegment of $\bigtriangleup ABG$ . Thus, the ratio of the area of $\bigtriangleup ADF$ to the area of $\bigtriangleup ABG$ is $1:4$, so the area of $\bigtriangleup ACG$ is $\frac{1}{4}\times100=25$. Therefore, the area of quadrilateral $FDBG$ is $[ABG]-[ADF]=100-25=\boxed{\textbf{(D) }75}$




\textbf{Solution 4}

The area of quadrilateral $FDBG$ is the area of $\bigtriangleup ABG$ minus the area of $\bigtriangleup ADF$. Notice, $\overline{DE} || \overline{BC}$, so $\bigtriangleup ABG \sim \bigtriangleup ADF$, and since $\overline{AD}:\overline{AB}=1:2$, the area of $\bigtriangleup ADF:\bigtriangleup ABG=(1:2)^2=1:4$. Given that the area of $\bigtriangleup ABC$ is $120$, using $\frac{bh}{2}$ on side $AB$ yields $\frac{50h}{2}=120\implies h=\frac{240}{50}=\frac{24}{5}$. Using the Angle Bisector Theorem, $\overline{BG}:\overline{BC}=50:(10+50)=5:6$, so the height of $\bigtriangleup ABG: \bigtriangleup ACB=5:6$. Therefore our answer is $\big[ FDBG\big] = \big[ABG\big]-\big[ ADF\big] = \big[ ABG\big]\big(1-\frac{1}{4}\big)=\frac{3}{4}\cdot \frac{bh}{2}=\frac{3}{8}\cdot 50\cdot \frac{5}{6}\cdot \frac{24}{5}=\frac{3}{8}\cdot 200=\boxed{\textbf{(D) }75}$




\textbf{Solution 5 (Trigonometry)}

We try to find the area of quadrilateral $FDBG$ by subtracting the area outside the quadrilateral but inside triangle $ABC$. Note that the area of $\triangle ADE$ is equal  to $\frac{1}{2} \cdot 25 \cdot 5 \cdot \sin{A}$ and the area of triangle $ABC$ is equal to $\frac{1}{2} \cdot 50 \cdot 10 \cdot \sin A$. The ratio $\frac{[ADE]}{[ABC]}$ is thus equal to $\frac{1}{4}$ and the area of triangle $ADE$ is $\frac{1}{4} \cdot 120 = 30$. Let side $BC$ be equal to $6x$, then $BG = 5x, GC = x$ by the angle bisector theorem. Similarly, we find the area of triangle $AGC$ to be $\frac{1}{2} \cdot 10 \cdot x \cdot \sin C$ and the area of triangle $ABC$ to be $\frac{1}{2} \cdot 6x \cdot 10 \cdot \sin C$. A ratio between these two triangles yields $\frac{[ACG]}{[ABC]} = \frac{x}{6x} = \frac{1}{6}$, so $[AGC] = 20$. Now we just need to find the area of triangle $AFE$ and subtract it from the combined areas of $[ADE]$ and $[ACG]$, since we count it twice. Note that the angle bisector theorem also applies for $\triangle ADE$ and $\frac{AE}{AD} = \frac{1}{5}$, so thus $\frac{EF}{ED} = \frac{1}{6}$ and we find $[AFE] = \frac{1}{6} \cdot 30 = 5$, and the area outside $FDBG$ must be $[ADE] + [AGC] - [AFE] = 30 + 20 - 5 = 45$, and we finally find  $[FDBG] = [ABC] - 45 = 120 -45 = \boxed{\textbf{(D) }75}$, and we are done.

 



\textbf{Solution 7 (Barycentrics)}

Let our reference triangle be $\triangle ABC$. Consequently, we have $A=(1,0,0)$, $B=(0,1,0)$, $C=(0,0,1).$ Since $D$ is the midpoint of $\overline{AB}$, we have that $D=(1:1:0)$. Similarly, we have $E=(1:0:1).$ Hence, the line through $D$ and $E$ is given by the equation 

 
 \[0 =  \begin{vmatrix} x & y & z\\ 1 & 1 & 0\\ 1 & 0 & 1 \end{vmatrix}\]

 Additionally, since all points on $\overline{AG}$ are characterized by $(t:1:5)$, we may plug in for $x,y,z$ to get $t=6$. Thus, we have $F=(6:1:5).$ Now, we homogenize the coordinates for $F D, B, G$ to get $F=(\frac{1}{2}, \frac{5}{12}, \frac{1}{12})$, $D=(\frac{1}{2}, \frac{1}{2}, 0)$, $B=(0,1,0)$, $G=(0, \frac{1}{6}, \frac{5}{6})$

 Splitting $[FBGD]$ into $[ DBG ] + [ FDG],$ we may now evaluate the two determinants:

 
 \[\begin{vmatrix} \frac{1}{2} & \frac{1}{2} & 0\\ 0 & 1 & 0\\ 0 & \frac{1}{6} & \frac{5}{6} \end{vmatrix}\]
 \[\begin{vmatrix} \frac{1}{2} & \frac{1}{12} & \frac{5}{12}\\ \frac{1}{2} & \frac{1}{2} & 0\\ 0 & \frac{5}{6} & \frac{1}{6} \end{vmatrix}.\]

 After simplification, we get $\frac{5}{12}$ and $\frac{5}{24}$, respectively. Summing, we get $\frac{15}{24}.$ Hence, $[FBGD]=\frac{15}{24} \cdot 120 = \fbox{\textbf{(D) }75}.$$\sim$Math0323

 



\textbf{Solution 8}


                    \begin{figure}[H]
                        \centering
                        \includegraphics[width=0.5\textwidth]{Solution/AMC_12_12A_2018_18_Solution_Figure_1.png}
                    \end{figure}
                    

 We want to find ratios. Start off letting [AEF]=1. Then we know [AEF]=[CEF]=1 since AE=CE and they share the same height. Then, also because AE is half of AC, we see that [AEF]=1/4 [ACG] so [ACG]=4. 

 But we also know [EFGC]=3/4 [ACG], so [EFGC]=3 and then [CFG]=2. If [CFG]=2, then note that [FGB] must be 2*5=10.

 Now, we see that [AFD]=5[AED] because EF=1/5 FD (from angle bisector theorem) and they share the same height. Then, because AD=DB, we have [FDB]=[FAD]=5. 

 Adding these up, we see that the ratio of [BGFD] to [ABC] is 15/24=5/8. 5/8*120 yields our answer of 75. ~mathboy282

 



\textbf{Solution 9}

Draw a nice diagram. Lets find some some areas. We will start by finding $[AED]$. The side length ratio of $\triangle{AED}$ to $\triangle{ACB}$ is $\dfrac{1}{2}$, so the area ratio is $\dfrac{1}{4}$, so $[AED] = \dfrac{1}{4} \cdot 120 = 30$. Lets look for information we have used yet. We still havent used the fact that $AG$ is an angle bisector. This hints at us using the angle bisector theorem. Letting $EF = y$, we have $FD = 5y$, and letting $CG = x$, $GB = 5x$. Now, we have
 \[\dfrac{[AEF]}{[AED]} = \dfrac{1}{6} \Longrightarrow [AEF] = \dfrac{1}{6} \cdot 30 = 5, [AFD] = \dfrac{5}{6} \cdot 30 = 25\]. Doing the same thing with triangles $ACG$ and $AGB$, we find that $[AGB] = 100$, and so $[FDGB] = 100-25 = \boxed{75}$

 -jb2015007

 


	\item (\textit{AMC 12 12A 2018}) Let $A$ be the set of positive integers that have no prime factors other than $2$, $3$, or $5$. The infinite sum 
 \[\frac{1}{1} + \frac{1}{2} + \frac{1}{3} + \frac{1}{4} + \frac{1}{5} + \frac{1}{6} + \frac{1}{8} + \frac{1}{9} + \frac{1}{10} + \frac{1}{12} + \frac{1}{15} + \frac{1}{16} + \frac{1}{18} + \frac{1}{20} + \cdots\]of the reciprocals of the elements of $A$ can be expressed as $\frac{m}{n}$, where $m$ and $n$ are relatively prime positive integers. What is $m+n$?


 $\textbf{(A) } 16 \qquad \textbf{(B) } 17 \qquad \textbf{(C) } 19 \qquad \textbf{(D) } 23 \qquad \textbf{(E) } 36$

 \textbf{Answer: }C



\textbf{Solution 1}

Note that the fractions of the form $\frac{1}{2^a3^b5^c},$ where $a,b,$ and $c$ are nonnegative integers, span all terms of the infinite sum.

 Therefore, the infinite sum becomes
 \begin{align*} \sum_{a=0}^{\infty}\sum_{b=0}^{\infty}\sum_{c=0}^{\infty}\frac{1}{2^a3^b5^c} &= \left(\sum_{a=0}^{\infty}\frac{1}{2^a}\right)\cdot\left(\sum_{b=0}^{\infty}\frac{1}{3^b}\right)\cdot\left(\sum_{c=0}^{\infty}\frac{1}{5^c}\right) \\ &= \frac{1}{1-\frac12}\cdot\frac{1}{1-\frac13}\cdot\frac{1}{1-\frac15} \\ &= 2\cdot\frac32\cdot\frac54 \\ &= \frac{15}{4} \end{align*}
by a product of geometric series, from which the answer is $15+4=\boxed{\textbf{(C) } 19}.$




\textbf{Solution 2}

This solution is the same as Solution 1 but potentially clearer.

 Clearly this is just summing over the reciprocals of the numbers of the form $2^i3^j5^k$, where $i,j,k\in [0,\infty)$. So our desired sum is $\sum_{k=0}^{\infty}\sum_{j=0}^{\infty}\sum_{i=0}^{\infty}\frac{1}{2^i3^j5^k}$. By the infinite geometric series formula, $\sum_{i=0}^{\infty}\frac{1}{2^i3^j5^k}$ is just $\frac{\frac{1}{3^j5^k}}{1-\frac{1}{2}}=\frac{2}{3^j5^k}$. Applying the infinite geometric series formula again gives that $\sum_{j=0}^{\infty}\frac{2}{3^j5^k}=\frac{\frac{2}{5^k}}{1-\frac{1}{3}}=\frac{3}{5^k}$. Applying the infinite geometric series formula again yields $\sum_{k=0}^{\infty}\frac{3}{5^k}=\frac{3}{1-\frac{1}{5}}=\frac{15}{4}$. Hence our final answer is $15+4=\boxed{\textbf{(C) } 19}$.

 



\textbf{Solution 3}

Separate into $7$ separate infinite series's so we can calculate each and find the original sum: 

 The first infinite sequence shall be all the reciprocals of the powers of $2$, the second shall be reciprocals of the powers of $3$, and the third will consist of reciprocals of the powers of $5$. We can easily calculate these to be $1, \frac{1}{2}, \frac{1}{4}$ respectively. 

 The fourth infinite series shall be all real numbers in the form $\frac{1}{2^a3^b}$, where $a,b\geq1$.

 The fifth is all real numbers in the form $\frac{1}{2^a5^b}$, where $a,b\geq1$. 

 The sixth is all real numbers in the form $\frac{1}{3^a5^b}$, where $a,b\geq1$. 

 The seventh infinite series is all real numbers in the form $\frac{1}{2^a3^b5^c}$, where $a,b,c\geq1$. 

 Let us denote the first sequence as $a_{1}$, the second as $a_{2}$, etc. We know $a_{1}=1$, $a_{2}=\frac{1}{2}$, $a_{3}=\frac{1}{4}$, let us find $a_{4}$. factoring out $\frac{1}{6}$ from the terms in this subsequence, we would get $a_{4}=\frac{1}{6}(1+a_{1}+a_{2}+a_{4})$. 

 Knowing $a_{1}$ and $a_{2}$, we can substitute and solve for $a_{4}$, and we get $\frac{1}{2}$. If we do similar procedures for the fifth and sixth sequences, we can solve for them too, and we get after solving them $\frac{1}{4}$ and $\frac{1}{8}$. 

 Finally, for the seventh sequence, we see $a_{7}=\frac{a_{8}}{30}$, where $a_{8}$  is the infinite series the problem is asking us to solve for. The sum of all seven subsequences will equal the one we are looking for, and we need to add the $\frac11$ term back: $1+1+\frac{1}{2}+\frac{1}{4}+\frac{1}{2}+\frac{1}{4}+\frac{1}{8}+\frac{a_{8}}{30}=a_{8}$. We solve this to get $\frac{29}{8}=\frac{29a_{8}}{30}$, or $\frac{15}{4}=a_{8}$. So our answer is $\frac{15}{4}$, but we are asked to add the numerator and denominator, which sums up to $\boxed{\textbf{(C) } 19}$.

 



\textbf{Solution 4(clear but takes time)}

First of all, we need to find some geometric sequences hidden in here.

 The first sequence is just $\frac{1}{1}$. We can just leave this 1 by itself.

 The second sequence contains all numbers in the form of $\frac{1}{2^p}$. This is just $\frac{1}{2} + \frac{1}{4} + \frac{1}{8} + ... = 1$.

 The third sequence contains all numbers in the form of $\frac{1}{2^p\cdot 3}$. These add to $\frac{\frac13}{\frac12}$.

 The fourth sequence contains all numbers in the form of $\frac{1}{2^p\cdot 5}$. These add to $\frac{\frac15}{\frac12}$.

 Essentially, every number is contained in a sequence that starts with the term that is the reciprocal of the result when the number is repeatedly divided by 2 until the number is odd.

 These will add up to $1+1+\frac{\frac{1}{3}+\frac{1}{5}+\frac{1}{9}+\frac{1}{15}+\frac{1}{25}+\frac{1}{27}+\cdots}{\frac{1}{2}}$. Note that the numerator of that fraction is the sum of the reciprocals of every number with only 3 or 5 in its prime factorization. Let's call that numerator $x$.

 So now, we find $x$. That strategy worked out pretty well. Let's use it again.

 The first sequence contains all numbers in the form of $\frac{1}{3^p}$. These add to $\frac{\frac13}{\frac23}$. The second sequence contains all numbers in the form of $\frac{1}{3^p\cdot5}$. These add to $\frac{\frac15}{\frac23}$. The third sequence contains all numbers in the form of $\frac{1}{3^p\cdot 5^2}$. These add to $\frac{\frac{1}{25}}{\frac{2}{3}}$

 Essentially, every number is contained in a sequence that starts with the term that is the reciprocal of the result when the number is repeatedly divided by 3 until the number is not divisible by 3.

 The total sum for $x$ is $\frac{1}{2}+\frac{\text{a geometric sequence}}{\frac{2}{3}}$, and that geometric sequence evaluates to $\frac{\frac{1}{5}}{1-\frac{1}{5}}=\frac{1}{4}$. Then we just need to do the computation.

 $x$ is $\frac{\frac14}{\frac23}+\frac12=\frac78$. Recalling our answer is $\frac{x}{\frac12}+2$, we can plug in $x$ to get a final answer of $\frac{7}{4}+2=\frac{15}{4}$. 

 Finally, we can add $15+4=\boxed{\textbf{(C) } 19}$.

 


	\item (\textit{AMC 12 12A 2018}) Triangle $ABC$ is an isosceles right triangle with $AB=AC=3$. Let $M$ be the midpoint of hypotenuse $\overline{BC}$. Points $I$ and $E$ lie on sides $\overline{AC}$ and $\overline{AB}$, respectively, so that $AI>AE$ and $AIME$ is a cyclic quadrilateral. Given that triangle $EMI$ has area $2$, the length $CI$ can be written as $\frac{a-\sqrt{b}}{c}$, where $a$, $b$, and $c$ are positive integers and $b$ is not divisible by the square of any prime. What is the value of $a+b+c$?


 $\textbf{(A) }9 \qquad \textbf{(B) }10 \qquad \textbf{(C) }11 \qquad \textbf{(D) }12 \qquad \textbf{(E) }13 \qquad$

 \textbf{Answer: }D



\textbf{Solution 1}

Observe that $\triangle{EMI}$ is isosceles right ($M$ is the midpoint of diameter arc $EI$ since $m\angle MEI = m\angle MAI = 45^\circ$), so $MI=2,MC=\frac{3}{\sqrt{2}}$. With $\angle{MCI}=45^\circ$, we can use Law of Cosines to determine that $CI=\frac{3\pm\sqrt{7}}{2}$. The same calculations hold for $BE$ also, and since $CI<BE$, we deduce that $CI$ is the smaller root, giving the answer of $\boxed{\textbf{(D) }12}$.

 



\textbf{Solution 2 (Ptolemy)}

We first claim that $\triangle{EMI}$ is isosceles and right.

 Proof: Construct $\overline{MF}\perp\overline{AB}$ and $\overline{MG}\perp\overline{AC}$. Since $\overline{AM}$ bisects $\angle{BAC}$, one can deduce that $MF=MG$. Then by AAS it is clear that $MI=ME$ and therefore $\triangle{EMI}$ is isosceles. Since quadrilateral $AIME$ is cyclic, one can deduce that $\angle{EMI}=90^\circ$. Q.E.D.

 Since the area of $\triangle{EMI}$ is 2, we can find that $MI=ME=2$, $EI=2\sqrt{2}$

 Since $M$ is the mid-point of $\overline{BC}$, it is clear that $AM=\frac{3\sqrt{2}}{2}$.

 Now let $AE=a$ and $AI=b$. By Ptolemy's Theorem, in cyclic quadrilateral $AIME$, we have $2a+2b=6$. By Pythagorean Theorem, we have $a^2+b^2=8$. One can solve the simultaneous system and find $b=\frac{3+\sqrt{7}}{2}$. Then by deducting the length of $\overline{AI}$ from 3 we get $CI=\frac{3-\sqrt{7}}{2}$, giving the answer of $\boxed{\textbf{(D) }12}$. (Surefire2019)

 



\textbf{Solution 3 (Elementary)}

Like above, notice that $\triangle{EMI}$ is isosceles and right, which means that $\dfrac{ME \cdot MI}{2} = 2$, so $MI^2=4$ and $MI = 2$. Then construct $\overline{MF}\perp\overline{AB}$ and $\overline{MG}\perp\overline{AC}$ as well as $\overline{MI}$. It's clear that $MG^2+GI^2 = MI^2$ by Pythagorean, so knowing that $MG = \dfrac{AB}{2} = \dfrac{3}{2}$ allows one to solve to get $GI = \dfrac{\sqrt{7}}{2}$. By just looking at the diagram, $CI=AC-MF-GI=\dfrac{3-\sqrt{7}}{2}$. The answer is thus $3+7+2=\boxed{\textbf{(D) }12}$.

 



\textbf{Solution 4 (Coordinate Geometry)}

Let $A$ lie on $(0,0)$, $E$ on $(0,y)$, $I$ on $(x,0)$, and $M$ on $\left(\frac{3}{2},\frac{3}{2}\right)$. Since ${AIME}$ is cyclic, $\angle EMI$ (which is opposite of another right angle) must be a right angle; therefore, $\overrightarrow{ME} \cdot \overrightarrow{MI} = \left<\frac{-3}{2}, y-\frac{3}{2}\right> \cdot \left<x-\frac{3}{2}, -\frac{3}{2}\right> = 0$. Compute the dot product to arrive at the relation $y=3-x$. We can set up another equation involving the area of $\triangle EMI$ using the \href{https://artofproblemsolving.com/wiki/index.php?title=Shoelace\_Theorem}{Shoelace Theorem}. This is 
 \[2=\frac{1}{2}\left[\frac{3}{2}\left(y-\frac{3}{2}\right)-xy+\frac{3}{2}\left(x+\frac{3}{2}\right)\right].\] Multiplying, substituting $3-x$ for $y$, and simplifying, we get $x^2 -3x + \frac{1}{2}=0$. Thus, $(x,y)=\left(\frac{3 \pm \sqrt{7}}{2},\frac{3 \mp \sqrt{7}}{2}\right)$. But $AI>AE$, meaning $x=AI=\frac{3 + \sqrt{7}}{2}$ and $CI = 3-\frac{3 + \sqrt{7}}{2}=\frac{3 - \sqrt{7}}{2}$, and the final answer is $3+7+2=\boxed{\textbf{(D) }12}$.

 



\textbf{Solution 5 (Quick)}

From $AIME$ cyclic we get $\angle{MEI} = \angle{MAI} = 45^\circ$ and $\angle{MIE} = \angle{MAE} = 45^\circ$, so $\triangle{EMI}$ is an isosceles right triangle. 

 From $[EMI]=2$ we get $EM=MI=2$. 

 Notice $\triangle{AEM} \cong \triangle{CIM}$, because $\angle{AEM}=180-\angle{AIM}=\angle{CIM}$, $EM=IM$, and $\angle{EAM}=\angle{ICM}=45^\circ$. 

 Let $CI=AE=x$, so $AI=3-x$. 

 By Pythagoras on $\triangle{EAI}$ we have $x^2+(3-x)^2=EI^2=8$, and solve this to get $x=CI=\dfrac{3-\sqrt{7}}{2}$ for a final answer of $3+7+2=\boxed{\textbf{(D) }12}$.

 



\textbf{Solution 6 (Bash)}

Let $CI=a$, $BE=b$. Because opposite angles in a cyclic quadrilateral are supplementary, we have $\angle EMI=90^{\circ}$. By the law of cosines, we have $MI^2=a^2+\frac{9}{4}-\frac{3}{2}a$, and $ME^2=b^2+\frac{9}{4}-\frac{3}{2}b$. Notice that $EI=2MO$, where $O$ is the origin of the circle mentioned in the problem. Thus $\frac{2MO*MO}{2}=2\implies MO=\sqrt{2}, EI=2\sqrt{2}$. By the Pythagorean Theorem, we have $ME^2+MI^2=EI^2\implies a^2+\frac{9}{4}-\frac{3}{2}a+b^2+\frac{9}{4}-\frac{3}{2}b=(2\sqrt{2})^2=8$. By the Pythagorean Theorem, we have $AE^2+AI^2=EI^2\implies (3-a)^2+(3-b)^2=(2\sqrt{2})^2=8\implies 18-6a-6b+a^2+b^2=8$. Thus we have $18-6a-6b+a^2+b^2=a^2+\frac{9}{4}-\frac{3}{2}a+b^2+\frac{9}{4}-\frac{3}{2}b\implies 18-6a-6b=\frac{9}{2}-\frac{3}{2}a-\frac{3}{2}b\implies \frac{27}{2}$ $=\frac{9}{2}a+\frac{9}{2}b\implies a+b=3\implies 3-b=a$. We know that
 \begin{align*} (3-a)^2+(3-b)^2&=8 \\ (3-a)^2+a^2&=8 \\ 2a^2-6a+9&=8 \\ 2a^2-6a+1&=0 \\ a&=\frac{6\pm \sqrt{36-8}}{2}=\frac{3\pm\sqrt{7}}{2}. \end{align*}
We take the smaller solution because we have $AI>AE\implies 3-AI<3-AE\implies CI<CE$, and we want $CI$, not $CE$, thus $CI=\frac{3-\sqrt{7}}{2}$. Thus our final answer is $3+7+2=\boxed{\textbf{(D) }12}$.

 -vsamc

 


	\item (\textit{AMC 12 12A 2018}) Which of the following polynomials has the greatest real root?


 $\textbf{(A) }   x^{19}+2018x^{11}+1   \qquad        \textbf{(B) }   x^{17}+2018x^{11}+1   \qquad    \textbf{(C) }   x^{19}+2018x^{13}+1   \qquad   \textbf{(D) }  x^{17}+2018x^{13}+1 \qquad  \textbf{(E) }   2019x+2018$

 \textbf{Answer: }B



\textbf{Solution 1 (Intermediate Value Theorem, Inequalities, Graphs)}

Denote the polynomials in the answer choices by $A(x),B(x),C(x),D(x),$ and $E(x),$ respectively.

 Note that $A(x),B(x),C(x),D(x),$ and $E(x)$ are strictly increasing functions with range $(-\infty,\infty).$ So, each polynomial has exactly one real root. The real root of $E(x)$ is $x=-\frac{2018}{2019}\approx-1.00.$ On the other hand, since $A(-1)=B(-1)=C(-1)=D(-1)=-2018$ and $A(0)=B(0)=C(0)=D(0)=1,$ we conclude that the real root for each of $A(x),B(x),C(x),$ and $D(x)$ must satisfy $x\in(-1,0)$ by the Intermediate Value Theorem (IVT).

 We analyze the polynomials for $x\in(-1,0):$

            \begin{enumerate} 
        \item We have
 \begin{align*} B(x)-A(x)=D(x)-C(x)&=x^{17}-x^{19} \\ &=x^{17}\left(1-x^2\right) \\ &<0. \end{align*}
As the graph of $y=A(x)$ is always above the graph of $y=B(x)$ in this interval, we deduce that $B(x)$ has a greater real root than $A(x)$ does. By the same reasoning, $D(x)$ has a greater real root than $C(x)$ does.

        \item We have
 \begin{align*} B(x)-D(x)&=2018x^{11}-2018x^{13} \\ &=2018x^{11}\left(1-x^2\right) \\ &<0, \end{align*}
from which $B(x)$ has a greater real root than $D(x)$ does.

            \end{enumerate}
            Now, we are left with comparing the real roots of $B(x)$ and $E(x).$ Since $B\left(-\frac{1}{\sqrt2}\right)<0<B(0),$ it follows that the real root of $B(x)$ must satisfy $x\in\left(-\frac{1}{\sqrt2},0\right)$ by the IVT. Clearly, $\boxed{\textbf{(B) }   x^{17}+2018x^{11}+1}$ has the greatest real root.

 



\textbf{Solution 2 (Similar to Solution 1)}

We can see that our real solution has to lie in the open interval $(-1,0)$. From there, note that $x^a < x^b$ if $a$, $b$ are odd positive integers if $a<b$, so hence it can only either be $\textbf{(B)}$ or $\textbf{(E)}$ (as all of the other polynomials will be larger than the polynomial $\textbf{(B)}$). Observe that $\textbf{(E)}$ gives the solution $x=-\frac{2018}{2019}$. We can approximate the root for $\textbf{(B)}$ by using $x=-\frac 12$: 
 \[\left(-\frac{1}{2}\right)^{17} - \frac{2018}{2048} + 1 \approx 0.\] Therefore, the root for $\textbf{(B)}$ is approximately $-\frac 12$. The answer is $\boxed{\textbf{(B) }   x^{17}+2018x^{11}+1}$.

 



\textbf{Solution 3 (Simpler Inequalities)}

The real roots must lie in the interval $(-1, 0).$ The root of choice E is around $-1$ which is the smallest value on this interval, so we can safely eliminate this. 

 Let's pick a number between $0$ and $-1,$ say $-1/2.$ We'll plug this in for $x.$ Note that plugging in $-1/2$ here will give us positive values for all polynomials. This means their roots are less than $-1/2$ since the graph increases from $x = -1$ to $0.$ Therefore, the polynomials giving the smallest values when $-1/2$ is plugged in for $x$ have roots closer to $-1/2$ (closer to the positive end of the interval) which means those roots are larger.

 
Comparing choices A and B:

 
 \[(-1/2)^{19} + 2018(-1/2)^{11} + 1 > (-1/2)^{17} + 2018(-1/2)^{11} + 1\]

 Comparing choices C and D:

 
 \[(-1/2)^{19} + 2018(-1/2)^{13} + 1 > (-1/2)^{17} + 2018(-1/2)^{13} + 1\]

 
We're left with choices B and D since they gave the smallest values:

 
 \[(-1/2)^{17} + 2018(-1/2)^{13} + 1 > (-1/2)^{17} + 2018(-1/2)^{11} + 1\]

 
Since choice B gave a smaller value than choice D, $\boxed{\textbf{(B) } x^{17}+2018x^{11}+1}$ has the largest roots.

 Desmos graph (you’ll need to zoom in): \href{https://www.desmos.com/calculator/bgibeqbiku}{https://www.desmos.com/calculator/bgibeqbiku}

 



\textbf{Solution 4 (Similar to Solution 1)}

Let the real solution to $B$ be $a.$ It is easy to see that when $a$ is plugged in to $A,$ since $-1 < a < 0,$ it follows that $a^{19} < a^{17}$ thus making the real solution to $A$ more "negative", or smaller than $B.$ Similarly we can assert that $D > C.$ Now to compare $B$ and $D,$ we can use the same method to what we used before to compare $B$ to $A,$ in which it is easy to see that the smaller exponent $(11)$ "wins". Now, the only thing left is for us to compare $B$ and $E.$ Plugging $\frac{-2018}{2019}$ (or the solution to $E$) into $B$ we obtain $\frac{(-2018)^{17}}{2019^{17}} + 2018\cdot\frac{(-2018)^{11}}{2019^{11}} + 1,$ which is intuitively close to $-1 - 2018 + 1 = -2018,$ much smaller than the solution the required $0.$ (For a more rigorous proof, one can note that $\left(\frac{2018}{2019}\right)^{17}$ and $\left(\frac{2018}{2019}\right)^{11}$ are both much greater than $\left(\frac{2018}{2019}\right)^{2019} \approx \frac{1}{e},$ by the limit definition of $e.$ Since $- \frac{1}{e} - 2018 \cdot \frac{1}{e} + 1$ is still much smaller the required $0$ for the solution to $B$ to be a solution, our answer is $\boxed{\textbf{(B) }   x^{17}+2018x^{11}+1}.$




\textbf{Solution 5}

Denote the polynomials in the answer choices by $A(x),B(x),C(x),D(x),$ and $E(x),$ respectively.

 The real root of $E(x)$ is $x=-\frac{2018}{2019}\approx-1.000.$, therefore, we can eliminate this polynomial first. Notice that all of the real roots for all answer choices are in $(-1,0)$. Therefore,
 \[x^{11}<0, \quad x^2<1, \quad x^{13}>x^{11}, \quad x^{19}+2018x^{13}+1>x^{19}+2018x^{11}+1, \quad C(x)>A(x)\]
 \[\quad x^{17}+2018x^{13}+1>x^{17}+2018x^{11}+1, \quad D(x)>B(x)\]
 \[x^{17}<0, \quad x^{19}>x^{17}, \quad x^{19}+2018x^{11}+1>x^{17}+2018x^{11}+1, \quad A(x)>B(x)\]

 Now compare the real roots of $A(x)$ and $D(x)$.
 \[A(x) - D(x) = x^{19}+2018x^{11}+1 - x^{17}-2018x^{13}-1 = x^{17}(x^2-1) - 2018x^{11}(x^2-1) = x^{11}(x^2-1)(x^6-2018)\]
 \[\because x^{11}<0, \quad 0<x^2<1, \quad 0<x^6<1, \quad \therefore A(x) - D(x)<0, \quad A(x) < D(x)\]

 \[C(x)>D(x)>A(x)>B(x)\]

 As $B(x)$ has the smallest value for the same $x$, and $A(x),B(x),C(x),D(x),$ are all monotonically increasing functions, its real root must be the greatest, thus, the answer is $\boxed{\textbf{(B) }   x^{17}+2018x^{11}+1}$




\textbf{Solution 6 (Monotonic Function)}

Denote the polynomials in the answer choices by $A(x),B(x),C(x),D(x),$ and $E(x),$ respectively.

 Notice that all of the real roots for all answer choices are in $(-1,0)$. The real root of $E(x)$ is $x=-\frac{2018}{2019}\approx-1.000.$, therefore, we can eliminate this polynomial first.

 First compare the real roots of $A(x)$ and $C(x)$:

 Let the real root of $A(x)$ be $a$, and the real root of $C(x)$ be $b$.
 \[a^{19}+2018a^{11}+1 = b^{19}+2018b^{13}+1, \quad a^{19}+2018a^{11} = b^{19}+2018b^{13}\]
 \[\because a^2 < 1, \quad a^{11} < 0, \quad \therefore a^{13}>a^{11}\]
 \[a^{19}+2018a^{11} < a^{19}+2018a^{13}, \quad b^{19}+2018b^{13} < a^{19}+2018a^{13}\]Let $f(x) = x^{19}+2018x^{13}$. Notice that $f(x)$ is a monotonically increasing function. As $f(b)<f(a)$, $b<a$, $A(x)$ has a bigger real root than $C(x)$.

 Similarly, $B(x)$ has a bigger real root than $D(x)$.

 To determine the polynomial with the greatest real root, now we only need to compare $A(x)$ and $B(x)$.

 Let the real root of $A(x)$ be $a$, and the real root of $B(x)$ be $b$.
 \[a^{19}+2018a^{11}+1 = b^{17}+2018b^{11}+1, \quad a^{19}+2018a^{11} = b^{17}+2018b^{11}\]
 \[\because a^2 < 1, \quad a^{17} < 0, \quad \therefore a^{19}>a^{17}\]
 \[a^{17}+2018a^{11} < a^{19}+2018a^{11}, \quad a^{17}+2018a^{11} < b^{17}+2018b^{11}\]Let $f(x) = x^{17}+2018x^{11}$. Notice that $f(x)$ is a monotonically increasing function. As $f(a)<f(b)$, $a<b$, $B(x)$ has a bigger real root than $A(x)$.

 Thus, the answer is $\boxed{\textbf{(B) }  x^{17}+2018x^{11}+1}$




\textbf{Solution 7 (Proof by Contradiction)}

Denote the polynomials in the answer choices by $A(x),B(x),C(x),D(x),$ and $E(x),$ respectively.

 Notice that all of the real roots for all answer choices are in $(-1,0)$. The real root of $E(x)$ is $x=-\frac{2018}{2019}\approx-1.000.$, therefore, we can eliminate this polynomial first.

 First compare the real roots of $A(x)$ and $C(x)$:

 Let the real root of $A(x)$ be $a$, and the real root of $C(x)$ be $b$.
 \[a^{19}+2018a^{11}+1 = b^{19}+2018b^{13}+1\]
 \[a^{11}(a^8+2018) = b^{11}(b^8+2018b^2)\]
 \[\frac{ a^{11} }{ b^{11} } = \frac{ b^8+2018b^2 }{ a^8+2018 }\]As $b^2<1$, $b^8+2018b^2<b^8+2018$, 
 \[\frac{ a^{11} }{ b^{11} } = \frac{ b^8+2018b^2 }{ a^8+2018 } < \frac{ b^8+2018 }{ a^8+2018 }\]

 Assume that $-1<a<b<0$. Thus, 
 \[-1<a^{11}<b^{11}<0, \quad 1<\frac{ a^{11} }{ b^{11} }, \quad 1< \frac{ b^8+2018 }{ a^8+2018 }\]
$\because -1<a<b<0, \quad \therefore 1>a^8>b^8>0, \quad a^8 + 1 > b^8+1, \quad \therefore \frac{b^8+2018}{a^8+2018}<1$

 \[\text{ Contradicts with } 1<\frac{b^8+2018}{a^8+2018}\]
$\therefore$ assumption $-1<a<b<0$ is false, $-1<b<a<0$, $A(x)$ has a bigger real root than $C(x)$.

 Similarly, $B(x)$ has a bigger real root than $D(x)$.

 To determine the polynomial with the greatest real root, now we only need to compare $A(x)$ and $B(x)$.

 Let the real root of $A(x)$ be $a$, and the real root of $B(x)$ be $b$.
 \[a^{19}+2018a^{11}+1 = b^{17}+2018b^{11}+1\]
 \[a^{19} - b^{17} = 2018( b^{11} - a^{11} )\]

 Assume that $-1<b<a<0$. Thus, 
 \[-1<b^{11}<a^{11}<0, \quad b^{11}-a^{11}<0, \quad 2018( b^{11} - a^{11} ) < 0, \quad a^{19} - b^{17} < 0\]

 \[-1<b^{17}<a^{17}<0, \quad a^2<1, \quad a^{17}<a^{19}, \quad b^{17}<a^{19}, \quad a^{19} - b^{17} > 0\]

 \[\text{ Contradicts with }a^{19} - b^{17} < 0\]
$\therefore$ assumption $-1<b<a<0$ is false, $-1<a<b<0$, $B(x)$ has a bigger real root than $A(x)$.

 Thus, the answer is $\boxed{\textbf{(B) }   x^{17}+2018x^{11}+1}$




\textbf{Solution 8 (Calculus)}

Note that $a(-1)=b(-1)=c(-1)=d(-1) < 0$ and $a(0)=b(0)=c(0)=d(0) > 0$. Calculating the definite integral for each function in the interval $[-1,0]$, we see that $B(x)\rvert^{0}_{-1}$ gives the most negative value. To maximize our real root, we want to maximize the area of the curve under the x-axis, which means we want our integral to be as negative as possible and thus the answer is $\boxed{\textbf{(B) }   x^{17}+2018x^{11}+1}$.

 



\textbf{Solution 9 (Calculus)}

Newton's Method is used to approximate the zero $x_{1}$ of any real valued function given an estimation for the root $x_{0}: \ x_{1}=x_{0}-{\frac {f(x_{0})}{f'(x_{0})}}.$ After looking at all the options, $x_{0}=-1$ gives a reasonable estimate. For options $\textbf{(A)}$ to $\textbf{(D)},$ we have $f(-1) = -2018$ and the estimation becomes $x_{1}=-1+{\frac {2018}{f'(-1)}}.$ Thus we need to minimize the derivative, giving us $\textbf{(B)}$. Now after comparing $\textbf{(B)}$ and $\textbf{(E)}$ through Newton's method, we see that $\textbf{(B)}$ has the higher root, so the answer is $\boxed{\textbf{(B) }   x^{17}+2018x^{11}+1}$.

 



\textbf{Solution 10 (very fast)}

Obviously, eliminate $(E).$ Note that for answer choices $(A)$ and $(B),$ $x \approx \sqrt[11]{\left(-\frac{1}{2018}\right)} > \sqrt[13]{\left(-\frac{1}{2018}\right)},$ with the RHS being approximately the value for $x$ when considering answer choices $(C)$ or $(D),$ so the answer is either $(A)$ or $(B).$ Now, since $x^{19} < x^{17}$ for all negative reals $x < -1$ (we want to choose the option closer to $0$ because then we will have to decrease $x$ less since we estimated), we have $\boxed{(B)}$ as the clear answer. 

 


	\item (\textit{AMC 12 12A 2018}) The solutions to the equations $z^2=4+4\sqrt{15}i$ and $z^2=2+2\sqrt 3i,$ where $i=\sqrt{-1},$ form the vertices of a parallelogram in the complex plane. The area of this parallelogram can be written in the form $p\sqrt q-r\sqrt s,$ where $p,$ $q,$ $r,$ and $s$ are positive integers and neither $q$ nor $s$ is divisible by the square of any prime number. What is $p+q+r+s?$
$\textbf{(A) } 20 \qquad  \textbf{(B) } 21 \qquad  \textbf{(C) } 22 \qquad  \textbf{(D) } 23 \qquad  \textbf{(E) } 24$

 \textbf{Answer: }A



\textbf{Solution 1 (Complex Numbers in Rectangular Form)}

We solve each equation separately:

 
            \begin{enumerate} 
        \item $z^2=4+4\sqrt{15}i$Let $z=a+bi$ for some real numbers $a$ and $b.$ 

 Substituting and expanding, we get 
 \begin{align*} (a+bi)^2&=4+4\sqrt{15}i \\ \left(a^2-b^2\right)+2abi&=4+4\sqrt{15}i. \end{align*}
Equating the real parts and the imaginary parts, respectively, we get
 \begin{align*} a^2-b^2&=4, &&(1) \\ ab&=2\sqrt{15}. &&(2) \end{align*}
We rearrange $(1)$ and square $(2):$
 \begin{align*} a^2&=b^2+4, \hspace{4mm} &&(1\star) \\ a^2b^2&=60. &&(2\star) \end{align*}
Substituting $(1\star)$ into $(2\star),$ we obtain $\left(b^2+4\right)b^2=60.$ Since $b^2\geq0$ for all real numbers $b,$ either inspection or factoring gives $b^2=6.$ Substituting this into either $(1\star)$ or $(2\star)$ produces $a^2=10.$ Since $ab>0$ from $(2),$ we have $(a,b)=\left(\sqrt{10},\sqrt{6}\right),\left(-\sqrt{10},-\sqrt{6}\right).$ 

 The solutions to $z^2=4+4\sqrt{15}i$ are $\boldsymbol{z=\sqrt{10}+\sqrt{6}i,-\sqrt{10}-\sqrt{6}i}.$


        \item $z^2=2+2\sqrt{3}i$By the same process, we have $(a,b)=\left(\sqrt3,1\right),\left(-\sqrt3,-1\right).$ 

 The solutions to $z^2=2+2\sqrt{3}i$ are $\boldsymbol{z=\sqrt3+i,-\sqrt3-i}.$

            \end{enumerate}
            Note that the problem is equivalent to finding the area of a parallelogram with consecutive vertices $(x_1,y_1)=\left(\sqrt{10}, \sqrt{6}\right),(x_2,y_2)=\left(\sqrt{3},1\right),(x_3,y_3)=\left(-\sqrt{10},-\sqrt{6}\right),$ and $(x_4,y_4)=\left(-\sqrt{3}, -1\right)$ in the coordinate plane. By the Shoelace Theorem, the area we seek is 
 \[\frac{1}{2} \left|(x_1y_2 + x_2y_3 + x_3y_4 + x_4y_1) - (y_1x_2 + y_2x_3 + y_3x_4 + y_4x_1)\right| = 6\sqrt2-2\sqrt{10},\] so the answer is $6+2+2+10=\boxed{\textbf{(A) } 20}.$




\textbf{Solution 2 (Complex Numbers in Polar Form)}

We solve each equation separately:

 
            \begin{enumerate} 
        \item $z^2=4+4\sqrt{15}i$Let $z=r(\cos\theta+i\sin\theta)=r\operatorname{cis}\theta,$ where $r$ is the magnitude of $z$ such that $r\geq0,$ and $\theta$ is the argument of $z$ such that $0\leq\theta<2\pi.$ 

 By De Moivre's Theorem, we have 
 \[z^2=r^2\operatorname{cis}(2\theta)=16\left(\frac14+\frac{\sqrt{15}}{4}i\right),\] from which

 
            \begin{itemize} <li>$r^2=16,$ so $r=4.$</li>
<li>$\begin{cases} \begin{aligned} \cos(2\theta) &= \frac14 \\ \sin(2\theta) &= \frac{\sqrt{15}}{4} \end{aligned}, \end{cases}$ so $\cos\theta=\pm\sqrt{\frac{1+\cos(2\theta)}{2}}=\pm\frac{\sqrt{10}}{4}$ and $\sin\theta=\pm\sqrt{\frac{1-\cos(2\theta)}{2}}=\pm\frac{\sqrt{6}}{4}$ by Half-Angle Formulas. 
Since $\cos(2\theta)>0$ and $\sin(2\theta)>0,$ it follows that $2\theta\in\biggl(0,\frac{\pi}{2}\biggr)\cup\biggl(2\pi,\frac{5\pi}{2}\biggr),$ or $\theta\in\biggl(0,\frac{\pi}{4}\biggr)\cup\biggl(\pi,\frac{5\pi}{4}\biggr).$ We conclude that $(r,\operatorname{cis}\theta)=\left(4,\frac{\sqrt{10}}{4}+\frac{\sqrt{6}}{4}i\right),\left(4,-\frac{\sqrt{10}}{4}-\frac{\sqrt{6}}{4}i\right).$ </li>

            \end{itemize}
            The solutions to $z^2=4+4\sqrt{15}i$ are $\boldsymbol{z=\sqrt{10}+\sqrt{6}i,-\sqrt{10}-\sqrt{6}i}.$

        \item $z^2=2+2\sqrt 3i,$By a similar process, we conclude that $(r,\theta)=\biggl(2,\frac{\pi}{6}\biggr),\biggl(2,\frac{7\pi}{6}\biggr).$ 

 The solutions to $z^2=2+2\sqrt{3}i$ are $\boldsymbol{z=\sqrt3+i,-\sqrt3-i}.$

            \end{enumerate}
            We continue with the last paragraph of Solution 1 to get the answer $\boxed{\textbf{(A) } 20}.$




\textbf{Solution 3 (Vectors)}

Let $z_1$ and $z_2$ be the solutions to the equation $z^2=4+4\sqrt{15}i,$ and $z_3$ and $z_4$ be the solutions to the equation $z^2=2+2\sqrt 3i.$ Clearly, $z_1$ and $z_2$ are opposite complex numbers, so are $z_3$ and $z_4.$ <i>\textbf{This solution refers to the results of De Moivre's Theorem in Solution 2.}</i>

 From Solution 2, let $z_1=4\operatorname{cis}\phi$ for some $0<\phi<\frac{\pi}{4}.$ It follows that $z_2=4\operatorname{cis}(\phi+\pi).$ On the other hand, we have $z_3=2\operatorname{cis}\frac{\pi}{6}$ and $z_4=2\operatorname{cis}\frac{7\pi}{6}$ without loss of generality. Since $\tan(2\phi)>\tan\frac{\pi}{3},$ we deduce that $2\phi>\frac{\pi}{3},$ from which $\phi>\frac{\pi}{6}.$

 In the complex plane, the positions of $z_1,z_2,z_3,$ and $z_4$ are shown below:
                    \begin{figure}[H]
                        \centering
                        \includegraphics[width=0.5\textwidth]{Solution/AMC_12_12A_2018_22_Solution_Figure_1.png}
                    \end{figure}
                    Note that the diagonals of every parallelogram partition the shape into four triangles with equal areas. Therefore, to find the area of the parallelogram with vertices $z_1,z_2,z_3,$ and $z_4,$ we find the area of the triangle with vertices $0,z_1,$ and $z_3,$ then multiply by $4.$

 Recall that $|z_1|=4, |z_2|=2, \sin\phi=\frac{\sqrt6}{4},$ and $\cos\phi=\frac{\sqrt{10}}{4}$ from Solution 2. The area of the parallelogram is
 \begin{align*} 4\cdot\left[\frac12\cdot|z_1|\cdot|z_3|\cdot\sin\left(\phi-\frac{\pi}{6}\right)\right] &= 16\sin\left(\phi-\frac{\pi}{6}\right) \\ &= 16\left[\sin\phi\cos\frac{\pi}{6}-\cos\phi\sin\frac{\pi}{6}\right] \\ &= 16\left[\frac{\sqrt3}{2}\sin\phi-\frac12\cos\phi\right] \\ &= 6\sqrt2-2\sqrt{10}, \end{align*}
so the answer is $6+2+2+10=\boxed{\textbf{(A) } 20}.$




\textbf{Solution 4 (Vectors)}

Rather than thinking about this with complex numbers, notice that if we take two solutions and think of them as vectors, the area of the parallelogram they form is half the desired area. Also, notice that the area of a parallelogram is $ab\sin \theta$ where $a$ and $b$ are the side lengths. 

 The side lengths are easily found since we are given the squares of $z$. Thus, the magnitude of $z$ in the first equation is just $\sqrt{16} = 4$ and in the second equation is just $\sqrt{4} = 2$. Now, we need $\sin \theta$. 

 To find $\theta$, think about what squaring is in complex numbers. The angle between the squares of the two solutions is twice the angle between the two solutions themselves. In addition, we can find $\cos$ of this angle by taking the dot product of those two complex numbers and dividing by their magnitudes. The vectors are $\Bigl\langle 4, 4\sqrt{15}\Bigr\rangle$ and $\Bigl\langle 2, 2\sqrt{3}\Bigr\rangle$, so their dot product is $8 + 24\sqrt{5}$. Dividing by the magnitudes yields: $\dfrac{8+24\sqrt{5}}{4 \cdot 16} = \dfrac{1 + 3\sqrt{5}}{8}$. This is $\cos 2\theta$, and recall the identity $\cos 2\theta = 1 - 2\sin^2 \theta$. This means that $\sin^2 \theta = \dfrac{7 - 3\sqrt{5}}{16}$, so $\sin \theta = \dfrac{\sqrt{7-  3\sqrt{5}}}{4}$. Now, notice that $\sqrt{7-  3\sqrt{5}} = \dfrac{3\sqrt{2}-\sqrt{10}}{2}$ (which is not too hard to discover) so $\sin \theta = \dfrac{3\sqrt{2}-\sqrt{10}}{8}$. Finally, putting everything together yields: $2\cdot 4 \cdot \dfrac{3\sqrt{2}-\sqrt{10}}{8} = 3\sqrt{2} - \sqrt{10}$ as the area of the parallelogram found by treating two of the solutions as vectors. However, drawing a picture out shows that we actually want twice this (each fourth of the parallelogram from the problem is one half of the parallelogram whose area was found above) so the desired area is actually $6\sqrt{2} - 2\sqrt{10}$. Then, the answer is $\boxed{\textbf{(A) } 20}$.

 



\textbf{Solution 5 (Shoelace, Similar to Solution 1)}

Let $z = x+yi$. Then, $z^2 = x^2-y^2 + 2xyi$. This means $x^2 - y^2 = 4$, $xy = 2\sqrt15$. We can quickly see $x=\sqrt{10}$ and $y=\sqrt{6}$. Another solution would be $x=-\sqrt{10}$, $y=-\sqrt{6}$.

 From the other equation, our solutions are $x= \pm \sqrt{3}$ and $y = \pm 1$.

 Then, our four coordinates are $A(\sqrt{10}, \sqrt{6})$, $B(-\sqrt{10}, -\sqrt{6})$, $C(\sqrt{3}, 1)$, $D(-\sqrt{3}, -1)$.

 Note that $ABCD$ isn't a convex parallelogram, so we'll use $ACBD$. Applying shoelace theorem, we get $6\sqrt{2} - 2\sqrt{10}$, which is $\boxed{\textbf{(A) } 20}$.

 


	\item (\textit{AMC 12 12A 2018}) In $\triangle PAT,$ $\angle P=36^{\circ},$ $\angle A=56^{\circ},$ and $PA=10.$ Points $U$ and $G$ lie on sides $\overline{TP}$ and $\overline{TA},$ respectively, so that $PU=AG=1.$ Let $M$ and $N$ be the midpoints of segments $\overline{PA}$ and $\overline{UG},$ respectively. What is the degree measure of the acute angle formed by lines $MN$ and $PA?$
$\textbf{(A) } 76 \qquad  \textbf{(B) } 77 \qquad  \textbf{(C) } 78 \qquad  \textbf{(D) } 79 \qquad  \textbf{(E) } 80$

 \textbf{Answer: }E



\textbf{Solution 1 (Trigonometry)}

Let $P$ be the origin, and $PA$ lie on the $x$-axis.

 We can find $U=\left(\cos(36), \sin(36)\right)$ and $G=\left(10-\cos(56), \sin(56)\right)$

 Then, we have $M=(5, 0)$ and $N$ is the midpoint of $U$ and $G$, or $\left(\frac{10+\cos(36)-\cos(56)}{2}, \frac{\sin(36)+\sin(56)}{2}\right)$

 Notice that the tangent of our desired angle is the absolute difference between the $y$-coordinates of the two points divided by the absolute difference between the $x$-coordinates of the two points.

 This evaluates to 
 \[\frac{\sin(36)+\sin(56)}{\cos(36)-\cos(56)}\]Now, using sum to product identities, we have this equal to 
 \[\frac{2\sin(46)\cos(10)}{-2\sin(46)\sin({-10})}=\frac{\sin(80)}{\cos(80)}=\tan(80)\]so the answer is $\boxed{\textbf{(E) } 80}.$

 Note: Though this solution is excellent, setting $M = (0,0)$ makes life a tad bit easier

 



\textbf{Solution 2 (Rotation, Isosceles Triangle, Parallel Lines)}

We will refer to the \textbf{Diagram} section. In this solution, all angle measures are in degrees.

 We rotate $\triangle PUM$ by $180^\circ$ about $M$ to obtain $\triangle AU'M.$ Let $H$ be the intersection of $\overline{PA}$ and $\overline{GU'},$ as shown below.
                    \begin{figure}[H]
                        \centering
                        \includegraphics[width=0.5\textwidth]{Solution/AMC_12_12A_2018_23_Solution_Figure_1.png}
                    \end{figure}
                    Note that $\triangle GU'A$ is an isosceles triangle with $GA=U'A=1,$ so $\angle AGU'=\angle AU'G=\frac{180-\angle GAU'}{2}=44.$ In $\triangle GHA,$ it follows that $\angle GHA=180-\angle GAH-\angle AGH=80.$

 Since $\frac{UM}{UU'}=\frac{UN}{UG}=\frac12,$ we conclude that $\triangle UMN\sim\triangle UU'G$ by SAS, from which $\angle UMN=\angle UU'G$ and $\angle UNM=\angle UGU'.$ By the Converse of the Corresponding Angles Postulate, we deduce that $\overline{MN}\parallel\overline{U'G}.$

 Finally, we have $\angle NMA=\angle GHA=\boxed{\textbf{(E) } 80}$ by the Corresponding Angles Postulate.

 



\textbf{Solution 3 (Extending PN)}

Link $PN$, extend $PN$ to $Q$ so that $QN=PN$. Then link $QG$ and $QA$. 

 $\because M,N$ are the midpoints of $PA$ and $PQ,$ respectively

 $\therefore MN$ is the midsegment of $\bigtriangleup PAQ$
$\therefore \angle QAP=\angle NMP$

 Notice that $\bigtriangleup PUN\cong \bigtriangleup QGN$

 As a result, $QG=AG=UP=1$, $\angle AQG=\angle QAG$, $\angle GQN=\angle NPU$

 Also, $\angle GQN+\angle QPA=\angle QPU+\angle QPA=\angle UPA=36^{\circ}$

 As a result, $2\angle QAG=180^{\circ}-56^{\circ}-36^{\circ}=88^{\circ}$

 Therefore, $\angle QAP=\angle QAG+\angle TAP=56^{\circ}+44^{\circ}=100^{\circ}$

 Since we are asked for the acute angle between the two lines, the answer to this problem is $\boxed{\textbf{(E) } 80}$.

 



\textbf{Solution 4 (Parallel Lines)}

Let the mid-point of $\overline{AT}$ be $B$ and the mid-point of $\overline{GT}$ be $C$. Since $BC=CG-BG$ and $CG=AB-\frac{1}{2}$, we can conclude that $BC=\frac{1}{2}$. Similarly, we can conclude that $BM-CN=\frac{1}{2}$. Construct $\overline{ND}\parallel\overline{BC}$ and intersects $\overline{BM}$ at $D$, which gives $MD=DN=\frac{1}{2}$. Since $\angle{ABD}=\angle{BDN}$, $MD=DN$, we can find the value of $\angle{DMN}$, which is equal to $\frac{1}{2}\angle T=44^{\circ}$. Since $\overline{BM}\parallel\overline{PT}$, which means $\angle{DMN}+\angle{NMP}+\angle{P}=180^{\circ}$, we can infer that $\angle{NMP}=100^{\circ}$. As we are required to give the acute angle formed, the final answer would be $80^{\circ}$, which is $\boxed{\textbf{(E) } 80}$.

 



\textbf{Solution 5 (Angle Bisectors)}

Let the bisector of $\angle ATP$ intersect $PA$ at $X.$ We have $\angle ATX = \angle PTX = 44^{\circ},$ so $\angle TXA = 80^{\circ}.$ We claim that $MN$ is parallel to this angle bisector, meaning that the acute angle formed by $MN$ and $PA$ is $80^{\circ},$ meaning that the answer is $\boxed{\textbf{(E) } 80}$.

 To prove this, let $N(x)$ be the midpoint of $U(x)G(x),$ where $U(x)$ and $G(x)$ are the points on $PT$ and $AT,$ respectively, such that $PU = AG = x.$ (The points given in this problem correspond to $x=1,$ but the idea we're getting at is that $x$ will ultimately not matter.) Since $U(x)$ and $G(x)$ vary linearly with $x,$ the locus of all points $N(x)$ must be a line. Notice that $N(0) = M,$ so $M$ lies on this line. Let $N(x_0)$ be the intersection of this line with $PT$ (we know that this line will intersect $PT$ and not $AT$ because $PT > AT$). Notice that $G(x_0) = T.$

 Let $AT = a, TP = b, PA = c.$ Then $AG(x_0) = PU(x_0) = AT = a$ and $PG(x_0) = PT = b.$ Thus, $PN(x_0) = \frac{a+b}{2}.$ By the Angle Bisector Theorem, $\frac{PX}{AX} = \frac{PT}{AT} = \frac{b}{a},$ so $PX = \frac{bc}{a+b}.$ Since $M$ is the midpoint of $AP,$ we also have $PM = \frac{c}{2}.$ Notice that:

 
 \[\frac{PM}{PX} = \frac{\frac{c}{2}}{\frac{bc}{a+b}} = \frac{a+b}{2b}\]
 \[\frac{PN(x_0)}{PT} = \frac{\frac{a+b}{2}}{b} = \frac{a+b}{2b}\]

 Since $\frac{PN(x_0)}{PT} = \frac{PM}{PX},$ the line containing all points $N(x)$ must be parallel to $TX.$ This concludes the proof.

 The critical insight to finding this solution is that the length $1$ probably shouldn't matter because a length ratio of $1:5$ or $1:10$ (as in the problem) is exceedingly unlikely to generate nice angles. This realization then motivates the idea of looking at all points similar to $N,$ which then leads to looking at the most convenient such point (in this case, the one that lies on $PT$).

 Shoutout to Richard Yi and Mark Kong for working with me to discover the necessary insights to this problem!

 



\textbf{Solution 6 (Overkill: Miquel Points)}

Note that $X$, the midpoint of major arc $PA$ on $(PAT)$ is the Miquel Point of $PUAG$ (Because $PU = AG$). Then, since $1 = \frac{UN}{NG} = \frac{PM}{MA}$, this spiral similarity carries $M$ to $N$. Thus, we have $\triangle XMN \sim \triangle XAG$, so $\angle XMN = \angle XAG$. 

 But, we have $\angle XAG = \angle PAG = \angle PAX = 56 - \frac{180 - \angle PXA}{2} =56 - \frac{180 - \angle T}{2} = 56 - \frac{\angle A + \angle P}{2} = 56 - \frac{56+36}{2} = 56 - 46 = 10$; thus $\angle XMN = 10$. 

 Then, as $X$ is the midpoint of the major arc, it lies on the perpendicular bisector of $PA$, so $\angle XMA = 90$. Since we want the acute angle, we have $\angle NMA = \angle XMA - \angle XMN = 90 - 10 = 80$, so the answer is $\boxed{\textbf{(E) } 80}$.

 <u>\textbf{Sidenote}</u>

 For another way to find $\angle XMN$, note that 
 \[\angle XAM = 90 - \angle MXA = 90 - \frac{\angle AXP}{2} = 90 - \frac{\angle ATP}{2}= 90 - 44 = 46,\] giving $\angle XMN = \angle XAG = 56 - 46 = 10$ as desired.

 



\textbf{Solution 7 (Olympiad Nuke)}

By \href{https://artofproblemsolving.com/community/c6h489748p2745891}{https://artofproblemsolving.com/community/c6h489748p2745891}, we get that $MN$ is parallel to the angle bisector of $\angle ATP.$ Thus, 
 \[\angle NMA = 180^\circ - 56^\circ - \frac{180^\circ - 56^\circ - 36^\circ}{2} = \boxed{\textbf{(E) } 80}.\]




\textbf{Solution 8 (Vectors)}

The argument of the average of any two unit vectors is average of the arguments of the two vectors. Thereby, the acute angle formed is 
 \[\frac{36^\circ{} + 180^\circ{} - 56^\circ{}}{2} = \boxed{\textbf{(E) } 80}.\]




\textbf{Solution 9}

Let $P$ be at point $(0, 0)$, $A$ be at point $(10, 0)$, and $M$ is at point $(5, 0)$. Let $U$ be at point $(x, y)$ and $G$ be at point $(a, b)$. We will use:

 $\tan(\theta) = |\frac{m_1 - m_2}{1 + m_1m_2}|$ where $\theta$ is the angle between two lines and $m_1, m_2$ are the slopes of the lines. Then we compute the slope of $PA$ which is simply $0$ because there is no change in $y$. The slope of $PU$ is $\frac{y}{x}$ and the slope of $AG$ is $\frac{b}{a - 10}$. We can plug in to get:

 $\tan(36) = |\frac{y}{x}|$ because thankfully the $0$ slope for $PA$ cancels a lot out. Clearly $\frac{y}{x}$ is positive so $\frac{y}{x} = \tan(36) \implies y = x\tan(36)$

 Similarly, $\tan(56) = |\frac{b}{a - 10}| = \frac{b}{10 - a}$ because $a < 10$ in the coordinate plane. Let $10 - a = c$. We have:

 $x^{2} + y^{2} = 1$ and $y = x\tan(36)$
$b^{2} + c^{2} = 1$ and $b = c\tan(56)$

 simply from substituting $10 - a = c$. 

 Note that $N$ is the midpoint of $UG$ so its coordinates are just: $N$ is at point $(\frac{x + a}{2}, \frac{y + b}{2})$ 

 We need to find the acute angle between $NM$ and $PA$ so again the $0$ slope cancels out a lot and we are left with:

 $\tan(\theta) = |\frac{y + b}{x + a - 10}| = |\frac{y + b}{x - c}|$

 All of these variables can be trivially solved from our equations above and we can substitute to get:

 $\tan(\theta) = |\frac{\sin(36) + \sin(56)}{\cos(36) - \cos(56)}|$

 Now we can use sum-to-product identities to simplify this. Note that $\sin(a) + \sin(b) = 2\sin(\frac{a + b}{2})\cos(\frac{a - b}{2})$ and $\cos(a) - \cos(b) = -2\sin(\frac{a + b}{2})\sin(\frac{a - b}{2})$ 

 Note that we can simplify using this to get $\tan(\theta) = |-\cot(\frac{36 - 56}{2})| = |-\cot(-10)| = \cot(10) = \tan(80)$ hence our answer is $\boxed{\textbf{(E)}}$



	\item (\textit{AMC 12 12A 2018}) Alice, Bob, and Carol play a game in which each of them chooses a real number between $0$ and $1.$ The winner of the game is the one whose number is between the numbers chosen by the other two players. Alice announces that she will choose her number uniformly at random from all the numbers between $0$ and $1,$ and Bob announces that he will choose his number uniformly at random from all the numbers between $\frac{1}{2}$ and $\frac{2}{3}.$ Armed with this information, what number should Carol choose to maximize her chance of winning?


 $\textbf{(A) }\frac{1}{2}\qquad \textbf{(B) }\frac{13}{24} \qquad \textbf{(C) }\frac{7}{12} \qquad \textbf{(D) }\frac{5}{8} \qquad \textbf{(E) }\frac{2}{3}\qquad$

 \textbf{Answer: }B



\textbf{Solution 1 (Expected Values)}

The expected value of Alice's number is $\frac12\left(0+1\right)=\frac12,$ and the expected value of Bob's number is $\frac12\left(\frac12+\frac23\right)=\frac{7}{12}.$ To maximize her chance of winning, Carol should choose the midpoint between these two expected values. So, the answer is $\frac12\left(\frac12+\frac{7}{12}\right)=\boxed{\textbf{(B) }\frac{13}{24}}.$

 Alternatively, once we recognize that the answer lies in the interval $\left(\frac12,\frac{7}{12}\right),$ we should choose $\textbf{(B)}$ since no other answer choices lie in this interval.

 



\textbf{Solution 2 (Piecewise Function)}

Let $a,b,$ and $c$ be the numbers that Alice, Bob, and Carol choose, respectively.

 Based on the value of $c,$ we construct the following table:
 \[\begin{array}{c|c|c} & & \\ [-2ex] \textbf{Case} & \textbf{Conditions for }\boldsymbol{a}\textbf{ and }\boldsymbol{b} & \textbf{Carol's Probability of Winning} \\ [0.5ex] \hline & & \\ [-1.5ex] 0<c<\frac12 & 0<a<c \text{ and } \frac12<b<\frac23 & \hspace{1.25mm}\frac{c}{1}\cdot\frac{1/6}{1/6}=c \\ [1.5ex] \frac12\leq c\leq\frac23 & \left(0<a<c \text{ and } c<b<\frac23\right) \text{ or } \left(c<a<1 \text{ and } \frac12<b<c\right) & \hspace{1.25mm}\frac{c}{1}\cdot\frac{2/3-c}{1/6}+\frac{1-c}{1}\cdot\frac{c-1/2}{1/6}=-12c^2+13c-3 \\ [1.5ex] \frac23<c<1 & c<a<1 \text{ and } \frac12<b<\frac23 & \hspace{4.375mm}\frac{1-c}{1}\cdot\frac{1/6}{1/6}=1-c \\ [1.5ex] \end{array}\]Let $P(c)$ be Carol's probability of winning when she chooses $c.$ We write $P(c)$ as a piecewise function:
 \[P(c) = \begin{cases} c & \mathrm{if} \ 0<c<\frac12 \\ -12c^2+13c-3 & \mathrm{if} \ \frac12\leq c\leq\frac23 \\ 1-c & \mathrm{if} \ \frac23<c<1 \end{cases}.\]Note that $P(c)$ is continuous in the interval $(0,1),$ increasing in the interval $\left(0,\frac12\right),$ increasing and then decreasing in the interval $\left(\frac12,\frac23\right),$ and decreasing in the interval $\left(\frac23,1\right).$ The graph of $y=P(c)$ is shown below.
                    \begin{figure}[H]
                        \centering
                        \includegraphics[width=0.5\textwidth]{Solution/AMC_12_12A_2018_24_Solution_Figure_1.png}
                    \end{figure}
                    Therefore, the maximum point of $P(c)$ occurs in the interval $\left[\frac12,\frac23\right],$ namely at $c=-\frac{13}{2\cdot(-12)}=\boxed{\textbf{(B) }\frac{13}{24}}.$




\textbf{Solution 3 (Answer Choices)}

Let $a,b,$ and $c$ be the numbers that Alice, Bob, and Carol choose, respectively.

 From the answer choices, we construct the following table:
 \[\begin{array}{c|c|c} & & \\ [-2ex] \boldsymbol{c} & \textbf{Conditions for }\boldsymbol{a}\textbf{ and }\boldsymbol{b} & \textbf{Carol's Probability of Winning} \\ [0.5ex] \hline & & \\ [-1.5ex] \frac12 & 0<a<\frac12 \text{ and } \frac12<b<\frac23 & \hspace{23.375mm}\frac{1/2}{1}\cdot\frac{1/6}{1/6}=\frac12 \\ [1.5ex] \frac{13}{24} & \left(0<a<\frac{13}{24} \text{ and } \frac{13}{24}<b<\frac23\right) \text{ or } \left(\frac{13}{24}<a<1 \text{ and } \frac12<b<\frac{13}{24}\right) & \frac{13/24}{1}\cdot\frac{1/8}{1/6}+\frac{11/24}{1}\cdot\frac{1/24}{1/6}=\frac{25}{48} \\ [1.5ex] \frac{7}{12} & \left(0<a<\frac{7}{12} \text{ and } \frac{7}{12}<b<\frac23\right) \text{ or } \left(\frac{7}{12}<a<1 \text{ and } \frac12<b<\frac{7}{12}\right) & \frac{7/12}{1}\cdot\frac{1/12}{1/6}+\frac{5/12}{1}\cdot\frac{1/12}{1/6}=\frac12 \\ [1.5ex] \frac58 & \left(0<a<\frac58 \text{ and } \frac58<b<\frac23\right) \text{ or } \left(\frac58<a<1 \text{ and } \frac12<b<\frac58\right) & \hspace{5.625mm}\frac{5/8}{1}\cdot\frac{1/24}{1/6}+\frac{3/8}{1}\cdot\frac{1/8}{1/6}=\frac{7}{16} \\ [1.5ex] \frac23 & \frac23<a<1 \text{ and } \frac12<b<\frac23 & \hspace{23.25mm}\frac{1/3}{1}\cdot\frac{1/6}{1/6}=\frac13 \\ [1.5ex] \end{array}\]Therefore, Carol should choose $\boxed{\textbf{(B) }\frac{13}{24}}$ to maximize her chance of winning.

 



\textbf{Solution 4 (Calculus)}

Note that Carol's number must lie in the interval $\left[\frac{1}{2}, \frac{2}{3}\right]$ because it never needs to be less than $\frac{1}{2}$ in order to be less than Bob's number, and it never needs to be greater than $\frac{2}{3}$ in order to be greater than Bob's number. Going past either value will only decrease the probability of being on the correct side of Alice's number.

 There are two cases of winning: 

 Case 1: Alice chooses a number that is smaller than Carol's, and Bob chooses a number that is bigger. 

 Case 2: Alice chooses a number that is bigger than Carol's, and Bob chooses a number that is smaller.

 Let Carol's number be $\frac{1}{2}+x$, where $x \in \left[0, \frac{1}{6}\right]$. The probability of Case 1 can be expressed as $\frac{\frac{1}{2} + x}{1}\cdot\frac{\frac{1}{6} - x}{\frac{1}{6}}=\left(\frac{1}{2} + x\right)\left(1 - 6x\right)$, and the probability of Case 2 can be expressed as $\frac{\frac{1}{2} - x}{1}\cdot\frac{x}{\frac{1}{6}}=\left(\frac{1}{2} - x\right)\left(6x\right)$. 

 Thus, the probability of Carol winning can be expressed as the sum of the probabilities of Cases 1 and 2: $P = \left(\frac{1}{2} + x\right)\left(1 - 6x\right) + \left(\frac{1}{2} - x\right)\left(6x\right)$, which simplifies to $P = \frac{1}{2} + x - 12x^2$. The maximum value of $P$ is obtained through the value of $x$ where the slope is $0$. We take the first derivative and get $1 - 24x$, which yields $0$ at $x = \frac{1}{24}$. Hence, Carol should select $\frac{1}{2} + \frac{1}{24} = \boxed{\textbf{(B) }\frac{13}{24}}$.

 Note that the same value of $x$ can be obtained through the Vertex Formula, $x=-\frac{b}{2a}$, without using Calculus.

 



\textbf{Solution 5 (Calculus)}

It suffices to find the average (expected) value of $C=\frac{\left(A+B\right)}{2}$ over the intervals $A \in \left[0,1\right]$ and $B \in \left[\frac{1}{2},\frac{2}{3}\right]$. We do this by finding $\int_0^1 \int_\frac{1}{2}^\frac{2}{3}\frac{\left(A+B\right)}{2}\,dB\,dA$ and divide by the area of the interval we're integrating over, namely ${\left(1-0\right)\left(\frac{2}{3}-\frac{1}{2}\right)}=\frac{1}{6}$. $\int_0^1 \left[\frac{AB}{2}+\frac{B^2}{4}\right]_{B=\frac{1}{2}}^\frac{2}{3}\,dA = \left[\frac{A^2}{24}+\frac{7A}{144}\right]_{A=0}^1=\frac{13}{144}$. Dividing by $\frac{1}{6}$ we get $\boxed{\textbf{(B) }\frac{13}{24}}$.

 



\textbf{Solution 6}

Draw a number line to visualize the situation. We have three important regions to consider: $0$ to $\frac{1}{2}$,  $\frac{1}{2}$ to $\frac{2}{3}$,  and  $\frac{2}{3}$ to $1$.

 If Carol chooses from $0$ to $\frac{1}{2}$, then Alice will be towards her left, and Bob will be towards her right. Carol should choose $\frac{1}{2}$ to maximize her chance of winning, since she wants to maximize her left for Alice to choose. In this case, the probability that Alice will be towards Carol's left is just $\frac{1}{2}$, so that is the probability of Carol winning. Meanwhile, if Carol chooses from $\frac{2}{3}$ to $1$, then Bob will be towards her left, and Alice will be towards her right. Carol should choose $\frac{2}{3}$ for the same reason, and her probability of winning is then $\frac{1}{3}$.

 At this point, it is clear that Carol should not choose the latter region, her probability of winning is less than the first.

 However, there is one more region to consider: $\frac{1}{2}$ to $\frac{2}{3}$. For this case, let Carol's number be C, and we will express the probability that Carol wins in terms of C.

 There are two cases.

 Case $1$: Alice towards her left, and Bob towards her right.

 The probability that Alice is towards Carol's left is simply $C/1$, or $C.$

 The probability that Bob is towards Carol's right is $(\frac{2}{3}-C)/(\frac{1}{6})$. 

 Therefore, the probability of Carol winning in this case is $C(4-6C)$.

 Case $2$: Bob towards her left, and Alice towards her right.

 Similarly, in this case, the probability of Carol winning is $(1-C)(6C-3)$.

 Adding, we find that the probability that Carol wins is $-12C^2 + 13C - 3.$

 To maximize this function, we either take the vertex value: $13/2(12) = 13/24$, or take the derivative and set it to zero.

 Either way, we find that the probability that Carol wins is $\boxed{\textbf{(B) }\frac{13}{24}}$, which is optimal since it is greater than $\frac{1}{2}$.

 


	\item (\textit{AMC 12 12A 2018}) For a positive integer $n$ and nonzero digits $a$, $b$, and $c$, let $A_n$ be the $n$-digit integer each of whose digits is equal to $a$; let $B_n$ be the $n$-digit integer each of whose digits is equal to $b$, and let $C_n$ be the $2n$-digit (not $n$-digit) integer each of whose digits is equal to $c$. What is the greatest possible value of $a + b + c$ for which there are at least two values of $n$ such that $C_n - B_n = A_n^2$?


 $\textbf{(A) } 12 \qquad \textbf{(B) } 14 \qquad \textbf{(C) } 16 \qquad \textbf{(D) } 18 \qquad \textbf{(E) } 20$

 \textbf{Answer: }D



\textbf{Solution 1}

By geometric series, we have
 \begin{alignat*}{8} A\_n&=a\bigl(\phantom{ }\underbrace{111\cdots1}\_{n\text{ digits}}\phantom{ }\bigr)&&=a\left(1+10+10\textasciicircum{}2+\cdots+10\textasciicircum{}{n-1}\right)&&=a\cdot\frac{10\textasciicircum{}n-1}{9}, \\ B\_n&=b\bigl(\phantom{ }\underbrace{111\cdots1}\_{n\text{ digits}}\phantom{ }\bigr)&&=b\left(1+10+10\textasciicircum{}2+\cdots+10\textasciicircum{}{n-1}\right)&&=b\cdot\frac{10\textasciicircum{}n-1}{9}, \\ C\_n&=c\bigl(\phantom{ }\underbrace{111\cdots1}\_{2n\text{ digits}}\phantom{ }\bigr)&&=c\left(1+10+10\textasciicircum{}2+\cdots+10\textasciicircum{}{2n-1}\right)&&=c\cdot\frac{10\textasciicircum{}{2n}-1}{9}. \end{alignat*}By substitution, we rewrite the given equation $C_n - B_n = A_n^2$ as
 \[c\cdot\frac{10^{2n}-1}{9} - b\cdot\frac{10^n-1}{9} = a^2\cdot\left(\frac{10^n-1}{9}\right)^2.\]Since $n > 0,$ it follows that $10^n > 1.$ We divide both sides by $\frac{10^n-1}{9}$ and then rearrange:
 \begin{align*} c\left(10^n+1\right) - b &= a^2\cdot\frac{10^n-1}{9} \\ 9c\left(10^n+1\right) - 9b &= a^2\left(10^n-1\right) \\ \left(9c-a^2\right)10^n &= 9b-9c-a^2. &&(\bigstar) \end{align*}
Let $y=10^n.$ Note that $(\bigstar)$ is a linear equation with $y,$ and $y$ is a one-to-one function of $n.$ Since $(\bigstar)$ has at least two solutions of $n,$ it has at least two solutions of $y.$ We conclude that $(\bigstar)$ must be an identity, so we get the following system of equations:
 \begin{align*} 9c-a^2&=0, \\ 9b-9c-a^2&=0. \end{align*}
The first equation implies that $c=\frac{a^2}{9}.$ Substituting this into the second equation gives $b=\frac{2a^2}{9}.$ 

 To maximize $a + b + c = a + \frac{a^2}{3},$ we need to maximize $a.$ Clearly, $a$ must be divisible by $3.$ The possibilities for $(a,b,c)$ are $(9,18,9),(6,8,4),$ or $(3,2,1),$ but $(9,18,9)$ is invalid. Therefore, the greatest possible value of $a + b + c$ is $6+8+4=\boxed{\textbf{(D) } 18}.$




\textbf{Solution 2}

Immediately start trying $n = 1$ and $n = 2$. These give the system of equations $11c - b = a^2$ and $1111c - 11b = (11a)^2$ (which simplifies to $101c - b = 11a^2$). These imply that $a^2 = 9c$, so the possible $(a, c)$ pairs are $(9, 9)$, $(6, 4)$, and $(3, 1)$. The first puts $b$ out of range but the second makes $b = 8$. We now know the answer is at least $6 + 8 + 4 = 18$.

 We now only need to know whether $a + b + c = 20$ might work for any larger $n$. We will always get equations like $100001c - b = 11111a^2$ where the $c$ coefficient is very close to being nine times the $a$ coefficient. Since the $b$ term will be quite insignificant, we know that once again $a^2$ must equal $9c$, and thus $a = 9, c = 9$ is our only hope to reach $20$. Substituting and dividing through by $9$, we will have something like $100001 - \frac{b}{9} = 99999$. No matter what $n$ really was, $b$ is out of range (and certainly isn't $2$ as we would have needed).

 The answer then is $\boxed{\textbf{(D) } 18}$.

 



\textbf{Solution 3}

The given equation can be written as
 \[c \cdot (\phantom{ } \overbrace{1111 \ldots 1111}^{2n\text{ digits}}\phantom{ }) - b \cdot (\phantom{ } \overbrace{11 \ldots 11}^{n\text{ digits}} \phantom{ }) = a^2 \cdot (\phantom{ } \overbrace{11 \ldots 11}^{n\text{ digits}} \phantom{ })^2.\]Divide by $\overbrace{11 \ldots 11}^{n\text{ digits}}$ on both sides:
 \[c \cdot (\phantom{ } \overbrace{1000 \ldots 0001}^{n+1\text{ digits}}\phantom{ }) - b = a^2 \cdot (\phantom{ } \overbrace{11 \ldots 11}^{n\text{ digits}} \phantom{ }).\]Next, split the first term to make it easier to deal with:
 \begin{align*} 2c + c \cdot (\phantom{ }\overbrace{99 \ldots 99}^{n\text{ digits}}\phantom{ }) - b &= a^2 \cdot (\phantom{ } \overbrace{11 \ldots 11}^{n\text{ digits}} \phantom{ }) \\ 2c - b &= (a^2 - 9c) \cdot (\phantom{ }\overbrace{11 \ldots 11}^{n\text{ digits}}\phantom{ }). \end{align*}
Because $2c - b$ and $a^2 - 9c$ are constants and because there must be at least two distinct values of $n$ that satisfy, $2c - b = a^2 - 9c = 0.$ Thus, we have
 \begin{align*} 2c&=b, \\ a^2&=9c. \end{align*}
Knowing that $a,b,$ and $c$ are single digit positive integers and that $9c$ must be a perfect square, the values of $(a,b,c)$ that satisfy both equations are $(3,2,1)$ and $(6,8,4).$ Finally, $6 + 8 + 4 = \boxed{\textbf{(D) } 18}.$




\textbf{Solution 4 (PaperMath's sum)}

By \href{https://artofproblemsolving.com/wiki/index.php?title=PaperMath\%E2\%80\%99s\_sum}{PaperMath’s sum}, the answer is at least $6+8+4=\boxed{\textbf{(D) } 18}.$



\end{enumerate}